\documentclass[11pt]{article}

\usepackage{fontspec}
\defaultfontfeatures{Renderer=Harfbuzz,Ligatures=TeX}
\directlua{luaotfload.add_fallback("clclatin", {"name:Noto Sans:mode=harf"})}
\directlua{luaotfload.add_fallback("clclatinmono", {"name:Noto Sans Mono:mode=harf"})}
\setmainfont[Script=Devanagari,RawFeature={fallback=clclatin}]{Noto Sans Devanagari}
\setsansfont[Script=Devanagari,RawFeature={fallback=clclatin}]{Noto Sans Devanagari}
\setmonofont[Script=Devanagari,RawFeature={fallback=clclatinmono}]{Noto Sans Devanagari}
\usepackage[margin=1in]{geometry}
\usepackage{amsmath,amssymb}
\usepackage{booktabs}
\usepackage{array}
\usepackage{longtable}
\usepackage{tabularx}
\usepackage{graphicx}
\usepackage{caption}
\usepackage{enumitem}
\usepackage{ragged2e}
\usepackage{hyperref}
\usepackage{url}
\usepackage{polyglossia}
\setdefaultlanguage{hindi}
\newfontfamily\devanagarifont[Script=Devanagari,RawFeature={fallback=clclatin}]{Noto Sans Devanagari}
\newfontfamily\devanagarifontsf[Script=Devanagari,RawFeature={fallback=clclatin}]{Noto Sans Devanagari}
\newfontfamily\devanagarifonttt[Script=Devanagari,RawFeature={fallback=clclatinmono}]{Noto Sans Devanagari}
\newfontfamily\latinfont{Noto Sans}
\newcolumntype{P}[1]{>{\raggedright\arraybackslash}p{#1}}
\captionsetup{font=small,labelfont=bf}
\setlength{\emergencystretch}{6em}
\hypersetup{
  unicode=true,
  colorlinks=true,
  linkcolor=blue,
  citecolor=blue,
  urlcolor=blue,
  pdflang={hi},
  pdftitle={Cosmo-Local Credit (CLC) White Paper v0.8},
  pdfauthor={William O. Ruddick and Mohamed Sohail}
}

\title{Cosmo-Local Credit (CLC)\\क्रेडिट को route करने, प्रतिबद्धताओं का समन्वय करने और स्वस्थ kosmo-local अर्थव्यवस्था को वित्त देने का नेटवर्क}
\author{William O. Ruddick \\ Mohamed Sohail \\ Grassroots Economics Foundation \\ \texttt{info@grassecon.org}}
\date{White Paper v0.8 translation: 10 October 2026}

\begin{document}
\maketitle
\tableofcontents
\newpage
\sloppy
\RaggedRight

\textbf{इस अनुवाद के बारे में।} यह White Paper v0.8 का अनुवाद है। मूल अंग्रेज़ी पाठ 30 सितंबर 2026 को प्रकाशित हुआ था। यह अनुवाद 10 अक्टूबर 2026 को प्रकाशित हुआ। अंग्रेज़ी मूल पाठ है। \href{https://docs.cosmolocal.credit/white-paper}{मूल अंग्रेज़ी पाठ पढ़ें}।

\textbf{संस्करण:} 0.8

\textbf{तारीख:} 30 सितंबर 2026

\textbf{लेखक:} William O. Ruddick और Mohamed Sohail

\textbf{संपर्क:} \href{mailto:info@grassecon.org}{info@grassecon.org}

\textbf{डाउनलोड:} \href{https://docs.cosmolocal.credit/white-paper/Cosmo-Local-Credit-CLC-White-Paper-v8-hi.pdf}{CLC White Paper v0.8 PDF हिंदी में}

\textbf{प्रकाशन:} अंतिम, संस्करणबद्ध सार्वजनिक प्रकाशन। यह दस्तावेज़ निवेश सलाह नहीं है और किसी प्रतिभूति या वित्तीय साधन को बेचने का प्रस्ताव या खरीदने का आग्रह नहीं है। बाद के संस्करण यहाँ बताए गए डिज़ाइनों में बदलाव कर सकते हैं।

\section*{इस श्वेत पत्र को कैसे पढ़ें}
\addcontentsline{toc}{section}{इस श्वेत पत्र को कैसे पढ़ें}

इस श्वेत पत्र में चार तरह की सामग्री है। यह स्थिति पूरे दस्तावेज़ पर लागू होती है, भले ही हर अध्याय इसे दोहराए नहीं।

\begin{longtable}{P{0.30\linewidth} P{0.62\linewidth}}
\toprule
स्थिति & अर्थ \\
\midrule
\endhead
\textbf{मौजूदा} & CLC App या सार्वजनिक Protocol v1.1.0 contracts में सत्यापित व्यवहार। इन contracts के लिए \href{https://docs.cosmolocal.credit/hi/protocol/overview}{प्रोटोकॉल संदर्भ} तथ्यात्मक स्रोत है। \\
\textbf{तैनाती पर निर्भर} & ऐसी क्षमता जो केवल तब मौजूद होती है जब कोई खास Pool, provider, अधिकार क्षेत्र, governing body या प्रकाशित policy उसे संभव बनाए। \\
\textbf{ऐतिहासिक} & पुरानी Sarafu Network सेवा का प्रमाण या कार्यक्षमता। यह CLC का मौजूदा उपयोग या यह प्रमाण नहीं कि कोई उपयोगकर्ता, संपत्ति, record या दायित्व स्थानांतरित हुआ। \\
\textbf{प्रस्तावित} & ऐसा design, economic model, governance mechanism या roadmap item जिसे लागू हुआ नहीं बताया गया है। \\
\bottomrule
\end{longtable}

मौजूदा Protocol v1.1.0 में सीधे \texttt{SwapPool} execution, वैकल्पिक चयन और मूल्यांकन, Pool token-balance caps, Pool fees, अतिरिक्त protocol fee और केवल quote देने वाला \texttt{SwapRouter} शामिल हैं। Multi-hop execution, HTLC या escrow routing, batch netting, साझा insurance, प्रस्तावित CLC governance token, प्रस्तावित CLC Network Pool, fee-credit mechanisms और network liquidity programs प्रस्तावित या तैनाती पर निर्भर हैं।

\textbf{किनके लिए:} Pool बनाने और चलाने वाले, protocol engineers और auditors, liquidity और mandate funders, governance designers तथा कानूनी या compliance reviewers।

\section*{सारांश}
\addcontentsline{toc}{section}{सारांश}

Cosmo-Local Credit एक उत्पाद परिवार है जो \textbf{Commitment Pooling Protocol (CPP)} के आसपास बना है। CPP बताता है कि स्वतंत्र रूप से संचालित प्रतिबद्धता पूल भुनाए जा सकने वाले वादों का चयन कैसे कर सकते हैं, विनिमय दर की विधियाँ और सीमाएँ कैसे प्रकाशित कर सकते हैं, भंडार कैसे रख सकते हैं और जवाबदेह आदान-प्रदान कैसे संभव बना सकते हैं।

\textbf{प्रतिबद्धता पूल } नियमों से संचालित पूरी व्यवस्था है। मौजूदा Protocol v1.1.0 का \textbf{\texttt{SwapPool}} उसका एक contract component है: token vault और सीधे swap का engine। तैनाती उस contract को registries, quoters, Pool token-balance caps, fee policies और दूसरी services से जोड़ सकती है।

Kosmo-local का अर्थ है कि खुले standards, software और knowledge को दुनिया भर में साझा किया जाए, जबकि issuing, वादा पूरा करना, governance और वास्तविक दुनिया की accountability स्थानीय रहें। जारीकर्ता अपने वाउचर के लिए ज़िम्मेदार रहता है। Pool प्रबंधक Pool के नियमों और स्पष्ट रूप से स्वीकार की गई किसी भी guarantee के लिए ज़िम्मेदार रहता है। न CLC App और न GEF अपने आप किसी voucher, Pool, value, route, liquidity position या outcome की guarantee देता है।

\section*{मौजूदा तैनाती}
\addcontentsline{toc}{section}{मौजूदा तैनाती}

GEF \href{https://cosmolocal.credit}{cosmolocal.credit} पर CLC App चलाता है। यह account और wallet access, सार्वजनिक Market catalog और tokens, vouchers, पेशकशों, प्रतिबद्धता पूलों, transfers और सीधे Pool swaps के लिए समर्थित interfaces देता है। उपलब्धता interface, contracts, inventory, configured limits और fees, network state, eligibility, providers, अधिकार क्षेत्र और जारीकर्ता या Pool की प्रकाशित शर्तों पर निर्भर करती है।

सिर्फ interface चलाने से GEF जारीकर्ता, Pool प्रबंधक, संरक्षक, ऋणदाता, उधारकर्ता, broker, guarantor, वाउचर भुनाने वाला, adviser, insurer या प्रतिभागियों के बीच दायित्वों का पक्ष नहीं बनता। App के इस्तेमाल पर \href{https://docs.cosmolocal.credit/hi/governance/terms}{सेवा की शर्तें} लागू होती हैं।

Protocol v1.1.0 जवाबदेह भूमिकाओं को तकनीकी नियंत्रण से अलग रखता है। Pool प्रबंधक, \texttt{SwapPool} मालिक, proxy administrator, dependency controller, catalog moderator और fee recipient अलग-अलग पक्ष हो सकते हैं। साझा शब्दावली के लिए \href{https://docs.cosmolocal.credit/hi/introduction/concepts}{अवधारणाएँ और शब्दावली} देखें।

\section*{ऐतिहासिक पृष्ठभूमि}
\addcontentsline{toc}{section}{ऐतिहासिक पृष्ठभूमि}

ऐतिहासिक Sarafu Network CLC App से पहले था और उसने सामुदायिक मुद्राओं और प्रतिबद्धताओं को एक साथ लाने के बारे में प्रमाण दिए। सेवा के Cosmo-Local Credit में बदलने से हर पुराना account, wallet, token, voucher, Pool, balance, report, participant या obligation स्थानांतरित नहीं हुआ।

इस श्वेत पत्र के ऐतिहासिक आँकड़े तारीख सहित एक internally consistent Sarafu dataset से आते हैं। ये CLC के मौजूदा उपयोग के आँकड़े नहीं हैं। स्रोत और माप की अवधि के लिए \href{https://docs.cosmolocal.credit/hi/white-paper/executive-summary}{कार्यकारी सारांश} देखें।

\section*{डिज़ाइन मॉडल}
\addcontentsline{toc}{section}{डिज़ाइन मॉडल}

CPP चार व्यापक कार्यों का इस्तेमाल करता है:

\begin{itemize}[leftmargin=*]
\item \textbf{चयन:} तय करना कि कौन से vouchers या दूसरी assets स्वीकार की जाएँ।
\item \textbf{मूल्यांकन:} Pool exchange rates या quotes बनाने की विधि प्रकाशित करना।
\item \textbf{सीमाएँ:} मौजूदा Pool token-balance caps या अलग से लागू दूसरी risk controls इस्तेमाल करना।
\item \textbf{आदान-प्रदान:} inventory रखना और बताई गई fees और सीमाओं के भीतर Pool swaps करना।
\end{itemize}
तैनाती routing, monitoring, liquidity support, guarantees, reserves, insurance या governance services केवल अलग से प्रकाशित policies के तहत जोड़ सकती है। Listing, चयन, कोई limit, reserve या blockchain record अपने आप जारीकर्ता की क्षमता, वादा पूरा होना, मंज़ूरी, guarantee या insurance साबित नहीं करता।

\section*{मूल्य और डिज़ाइन सिद्धांत}
\addcontentsline{toc}{section}{मूल्य और डिज़ाइन सिद्धांत}

\begin{itemize}[leftmargin=*]
\item \textbf{लोगों की देखभाल:} व्यक्तिगत और सामूहिक कल्याण का समर्थन।
\item \textbf{पर्यावरण की देखभाल:} पर्यावरणीय नुकसान कम करना और पुनर्योजी कार्यों का समर्थन।
\item \textbf{निष्पक्षता:} संसाधनों, ज्ञान और देखभाल तक सुरक्षित और न्यायसंगत पहुँच बढ़ाना।
\item \textbf{पारस्परिकता:} प्रकाशित नियमों के तहत जोखिम, लागत, ज़िम्मेदारी और अतिरिक्त लाभ साझा करना।
\item \textbf{वर्चस्व से बचाव:} लोगों, data, finance, knowledge और material resources पर केंद्रित नियंत्रण का विरोध।
\item \textbf{लचीलापन:} समुदायों को आर्थिक, राजनीतिक और climate disruptions के लिए तैयार होने और ढलने में मदद करना।
\item \textbf{स्थानीय जवाबदेही:} वास्तविक दुनिया में वादा पूरा करने और कानूनी ज़िम्मेदारी को पहचाने गए पक्षों के पास रखना।
\item \textbf{Forkability:} compatible communities और deployments को valid balances या obligations मिटाए बिना shared registries या services छोड़ने देना।
\end{itemize}
Digital layer साझा memory और coordination infrastructure है। यह relationships, जारीकर्ता के काम, local governance या legal accountability की जगह नहीं लेती।

\section*{मौजूदा और प्रस्तावित प्रक्रियाएँ}
\addcontentsline{toc}{section}{मौजूदा और प्रस्तावित प्रक्रियाएँ}

\subsection*{मौजूदा सीधी Pool प्रक्रिया}
\addcontentsline{toc}{subsection}{मौजूदा सीधी Pool प्रक्रिया}

\begin{enumerate}[leftmargin=*]
\item धारक voucher, उसके जारीकर्ता, उसकी शर्तों और उससे जुड़ी पेशकशों की जाँच करता है।
\item Pool समर्थित assets स्वीकार करता है और अपनी configuration और ज़िम्मेदार authorities प्रकाशित करता है।
\item धारक quote लेता है और सीधे Pool swap की अनुमति देता है।
\item \texttt{SwapPool} blockchain पर swap settlement करता है और Pool तथा protocol fees दर्ज करता है।
\item अगर धारक बाद में App में \textbf{`भुनाएँ'} चुनता है, तो App token owner को हस्तांतरण तैयार करता है ताकि voucher भुनाने के लिए पेश किया जा सके।
\item जारीकर्ता अलग से अपना वादा पूरा करता है और निस्तारण दर्ज करता है, ताकि पूरी की गई units दोबारा इस्तेमाल न हों।
\end{enumerate}
Pool से inventory लेना swap है। यह केवल तब voucher भुनाना है जब Pool ने अलग से जारीकर्ता की भूमिका स्वीकार की हो और लागू voucher terms पूरी करता हो।

\subsection*{व्यापक प्रस्तावित डिज़ाइन}
\addcontentsline{toc}{subsection}{व्यापक प्रस्तावित डिज़ाइन}

प्रस्तावित डिज़ाइन execution routing, netting, shared services, network liquidity programs, insurance policies और governance assets की संभावना देखता है। लागू होने से पहले हर व्यवस्था को implementation, ज़िम्मेदार पक्ष, स्वीकार किए गए rules, funding, controls और public disclosures चाहिए।

प्रस्तावित fee model Pool fees से ली गई network cut और अलग routing या service fees इस्तेमाल कर सकता है। यह प्रस्ताव Protocol v1.1.0 से अलग है, जहाँ वैकल्पिक protocol fee Pool fee के ऊपर लगती है।

मौजूदा \texttt{SwapPool} से समर्थकों या liquidity providers को अपने आप Pool shares, repayment, withdrawal, reward या governance rights नहीं मिलते। हर मौजूदा या भविष्य के program को transfer का प्रकार, ज़िम्मेदार पक्ष, control, recovery या withdrawal rights, losses, rewards, fees और governance rights अलग से बताने होंगे।

\section*{पढ़ने का मार्ग}
\addcontentsline{toc}{section}{पढ़ने का मार्ग}

\begin{itemize}[leftmargin=*]
\item मौजूदा शब्दावली और कार्रवाई का जीवनचक्र: \href{https://docs.cosmolocal.credit/hi/introduction/concepts}{अवधारणाएँ और शब्दावली}
\item सत्यापित और लागू contracts: \href{https://docs.cosmolocal.credit/hi/protocol/overview}{Protocol v1.1.0}
\item ऐतिहासिक प्रमाण: \href{https://docs.cosmolocal.credit/hi/white-paper/executive-summary}{कार्यकारी सारांश}
\item Confederation और interoperability: अध्याय 5, 8 और 11
\item प्रस्तावित guarantees और insurance की सीमाएँ: अध्याय 10 और 11
\item प्रस्तावित liquidity और fee models: अध्याय 7, 9 और 12
\item प्रस्तावित measurement framework: अध्याय 14 और परिशिष्ट A--C
\item प्रस्तावित roadmap: अध्याय 15
\item कानूनी और compliance विचार: अध्याय 17
\end{itemize}
\href{https://docs.cosmolocal.credit/hi/white-paper/archive}{श्वेत पत्र के पिछले संस्करण} केवल स्पष्ट रूप से बदले गए ऐतिहासिक प्रकाशनों के रूप में रखे गए हैं।

\section*{कार्यकारी सारांश}
\addcontentsline{toc}{section}{कार्यकारी सारांश}

वे \textbf{प्रतिबद्धता संवर्धन प्रोटोकॉल (CPP)} यह वर्णन करता है कि स्वतंत्र रूप से कैसे शासन किया जाता हैप्रतिबद्धता पूलमुआवजे के लिए प्रतिबद्धताएं संसाधित कर सकते हैं, विनिमय दर की विधियों और सीमाओं को प्रकाशित कर सकते है, इन्वेंट्री रख सकते हैं और जिम्मेदार आदान-प्रदान सक्षम कर सकते हैंः जारीकर्ता अपने वाउचर प्रतिबद्धताओं के लिए ज़िम्मेदार रहते हैं।पूल प्रबंधकइस तरह के नियम और उनके द्वारा स्पष्ट रूप से लिए जाने वाले किसी भी गारंटी के लिए जिम्मेदार हैं।CLC App और न हीGEF एक सार्वभौमिक गारंट है।

Protocol v1.1.0 वर्तमान प्रत्यक्ष स्वैप निर्माण ब्लॉक प्रदान करता हैः\texttt{GiftableToken}, \texttt{SwapPool}, वैकल्पिक रजिस्टर और मूल्यांकन मॉड्यूल, पूल टोकन-बालेंस कैप, पूल शुल्क, एक अतिरिक्त प्रोटोकॉल शुल्क, और केवल उद्धरण \texttt{SwapRouter}. वर्तमान अनुबंध वास्तविक दुनिया में वादे की पूर्ति दर्ज नहीं करते, स्वचालित तरलता प्रदाता शेयरों का निर्माण नहीं करते हैं, मल्टी-हॉप मार्गों को निष्पादित नहीं करते या सार्वभौमिक भंडार, गारंटी या बीमा प्रदान नहीं करते।

ऐतिहासिक \textbf{Sarafu Network} समुदाय मुद्राओं, बचत समूहों, पारस्परिक सहायता प्रणालियों और उत्पादन प्रतिबद्धताओं में प्रतिबद्धता साझा करने की अवधारणाओं का उपयोग किया। इसकी सेवा बाद में CLC App पर स्थानांतरित हो गई। यह एक प्रतिनिधित्व नहीं था कि हर ऐतिहासिक खाता, वॉलेट, टोकन, वाउचर, पूल, बैलेंस, रिपोर्ट, प्रतिभागी या दायित्व माइग्रेट हुआ।

\subsection*{ऐतिहासिक Sarafu Network स्नैपशॉट  5 जुलाई 2023 से 20 जुलाई 2025 तक सेलो गतिविधि}
\addcontentsline{toc}{subsection}{ऐतिहासिक Sarafu Network स्नैपशॉट  5 जुलाई 2023 से 20 जुलाई 2025 तक सेलो गतिविधि}

\textbf{स्रोत:} \href{https://dune.com/grassrootseconomics/sarafu-network}{ड्यून एनालिटिक्स डैशबोर्ड}

\begin{itemize}[leftmargin=*]
\item 26,367 उपयोगकर्ता
\item 285,197 सहकर्मी विनिमय
\item 188 अद्वितीय सक्रिय प्रतिबद्धता पूल
\item 745 अद्वितीय सक्रिय वाउचर
\item \$320,692 पूल स्वैप की मात्रा
\item 899 प्रकाशित रिपोर्ट (\href{https://sarafu.network/reports}{ऐतिहासिक रिपोर्ट अभिलेखागार})
\end{itemize}
ये आंकड़े ऐतिहासिक साक्ष्य हैं, वर्तमान CLC उपयोग के आंकड़े नहीं।

व्यापक प्रस्तावित CLC डिजाइन पता चलता है कि कैसे संगत नेटवर्क कर सकते हैंः

\begin{enumerate}[leftmargin=*]
\item स्वतंत्र रूप से शासित पूल के पार मार्गों की खोज और उद्धरण;
\item स्थानीय जिम्मेदारी के बिना जिम्मेदार नेटवर्क सेवाओं का समन्वय करना;
\item अलग से वित्त पोषित तरलता जनादेशों, गारंटीओं, भंडार या बीमा पॉलिसी का समर्थन करना;
\item पूल स्वैप को वाउचर प्रस्तुति, वादे की पूर्ति और डिस्चार्ज से अलग रूप से मापें; तथा
\item अपनाए गए साझा सेवाओं को वित्त पोषित करने के लिए एक प्रस्तावित नेटवर्क रैक और सेवा शुल्क का उपयोग करें।
\end{enumerate}
प्रस्ताव में एक \textbf{प्रस्तावित CLC शासन टोकन } और एक \textbf{ प्रस्तावित CLC नेटवर्क पूल}. न ही वर्तमान ऐप या Protocol v1.1.0 का हिस्सा है।

प्रस्तावित मॉडल के तहत, तरलता और शासन प्रतिभागियों को केवल प्रकाशित पात्रता, जोखिम, नियंत्रण, हस्तांतरण अधिकारों, वसूली की शर्तें और शुल्क नियमों के साथ अलग से अपनाए गए कार्यक्रमों के माध्यम से कार्य कर सकते थे। सामान्य पूल जमा स्वचालित शेयर, वापसी, निकासी, इनाम या शासन अधिकार नहीं बनाते हैं।

साझा उद्देश्य हैः

देखभाल, निष्पक्षता, स्थानीय जिम्मेदारी और लचीलापन को बनाए रखते हुए जिम्मेदार आदान-प्रदान और वास्तविक दुनिया की प्रतिबद्धताओं को पूरा करना बढ़ाना।

एंटी-कैप्चर और विश्वसनीय बाहर निकलने के मुख्य डिजाइन लक्ष्य बने हुए हैं। प्रस्तावित नियंत्रणों में समय में देरी वाले शासन, कई अनुमोदन सीमाएं, पारदर्शी प्रतिनिधिमंडल, संघर्ष नियम, घटना समीक्षा और एक प्रलेखित फोर्क-एंड-मिग्रेट प्रक्रिया शामिल है। ये नियंत्रण चुनिंदा शासन जोखिमों को कम करेंगे लेकिन नुकसान या गारंटी परिणामों को समाप्त नहीं कर पाएंगे।


\section*{1. प्रतिबद्धता संवर्धन प्रोटोकॉल (CPP): मूल आदिम}
\addcontentsline{toc}{section}{1. प्रतिबद्धता संवर्धन प्रोटोकॉल (CPP): मूल आदिम}

\textbf{मानसिक मॉडल:} प्रतिबद्धता पूल वाउचर या अन्य परिसंपत्तियों को स्वीकार करने, विनिमय नियमों को प्रकाशित करने, इन्वेंट्री रखने और स्वैप सक्षम करने के लिए एक नियंत्रित व्यवस्था है। धारक बाद में वास्तविक दुनिया में वादे की पूर्ति के लिए अपने जारीकर्ताओं को वाउचर्स प्रस्तुत करते हैं। पूल एक्सचेंज और जारीकर्ता वादे की पूर्ति अलग-अलग जीवन चक्र हैं।

CPP स्पष्ट रूप से वर्णित प्रतिबद्धताओं के माध्यम से मूल्य को समन्वयित करता है। \href{https://willruddick.substack.com/p/grassroots-economics-the-book-is}{जमीनी अर्थशास्त्रः चिंतन और अभ्यास}.

\subsection*{1.1 प्रतिबद्धता क्या है?}
\addcontentsline{toc}{subsection}{1.1 प्रतिबद्धता क्या है?}

एक प्रतिबद्धता किसी पहचाने गए पक्ष की भविष्य में वितरण का वादा है---उदाहरण के लिए भोजन, परिवहन, श्रम, भंडारण या अन्य वैध वस्तु, सेवा, लाभ या प्रदर्शन। \textbf{वाउचर} एक टोकन या रिकॉर्ड है जो प्रकाशित शर्तों के तहत उस प्रतिबद्धता के रूप में दर्शाया गया है।

टोकन अनुबंध डिजिटल यांत्रिकी को रिकॉर्ड करता है। वाउचर शर्तें जारीकर्ता, प्रस्ताव, क्षमता, स्थान, समय, प्रतिबंध, प्रस्तुति, वादे की पूर्ति, शिकायत और डिस्चार्ज प्रक्रिया की पहचान करती हैं।

\subsection*{1.2 प्रतिबद्धता पूल क्या है?}
\addcontentsline{toc}{subsection}{1.2 प्रतिबद्धता पूल क्या है?}

प्रतिबद्धता पूल एक नियंत्रित व्यवस्था है। इसे एक व्यक्ति, सहकारी, सामुदायिक समूह, सार्वजनिक एजेंसी, संघ, मल्टीसिग, सेवा ऑपरेटर या अन्य जिम्मेदार संरचना द्वारा प्रबंधित किया जा सकता है।

प्रासंगिक भूमिकाओं और अधिकारियों में निम्नलिखित शामिल हैंः

\begin{itemize}[leftmargin=*]
\item \textbf{पूल प्रबंधक:} पूल के नियमों और किसी भी स्पष्ट रूप से ली गई गारंटी को प्रकाशित करता है और प्रशासित करता है;
\item \textbf{पूल मालिकः} वर्तमान \texttt{SwapPool} मालिक अधिकारों को धारण करता है;
\item \textbf{प्रॉक्सी प्रशासक:} एक प्रॉक्सी कार्यान्वयन को अपग्रेड कर सकते हैं;
\item \textbf{निर्भरता नियंत्रक:} कॉन्फ़िगर किए गए रजिस्टर, कोटर्स, लिमिटर या शुल्क घटकों को नियंत्रित करें;
\item \textbf{मार्ग खोजक या संचालकः} उद्धरणों की खोज कर सकते हैं या, भविष्य के कार्यान्वयन में, एक अलग से अधिकृत मार्ग निष्पादित कर सकते है; तथा
\item \textbf{गारंटर:} केवल प्रकाशित, वित्त पोषित शर्तों के माध्यम से एक परिभाषित दायित्व लेता है।
\end{itemize}
CPP समूह पूल चार अवधारणाओं में कार्य करता हैः

\begin{itemize}[leftmargin=*]
\item \textbf{रखरखाव} समर्थित टोकन या वाउचर स्वीकार करें।
\item \textbf{मूल्यांकनः} एक विनिमय दर या बोली के लिए उपयोग की गई विधि प्रकाशित करें।
\item \textbf{सीमा:} वर्तमान तांडव टोकन बैलेंस कैप या अलग से लागू किए गए अन्य नियंत्रण लागू करें।
\item \textbf{विनिमयः} इन्वेंट्री रखना, स्वैप निष्पादित करना, शुल्क का लेखांकन करना और लेनदेन रिकॉर्ड जारी करना।
\end{itemize}
Protocol v1.1.0 इन कार्यों को लागू करता है\texttt{SwapPool} और वैकल्पिक निर्भरता।\texttt{Limiter} एक पूल में टोकन बैलेंस को कैप करता है; यह रोलिंग, प्रति खाता या नेटवर्क-व्यापी स्वैप सीमाएं प्रदान नहीं करता है।\texttt{SwapRouter} बहु पूल कोट्स की गणना करता है; यह स्वैप निष्पादित नहीं करता है।

व्यापक प्रस्तावितCPPडिजाइन रोलिंग सीमाओं, खाता नियंत्रण, निष्पादन राउटर जोड़ सकता है,HTLCवे प्रस्तावित घटक हैं, वर्तमान लिमिटर या राउटर के विवरण नहीं।

\subsection*{1.3 वर्तमान प्रत्यक्ष स्वैप तर्क}
\addcontentsline{toc}{subsection}{1.3 वर्तमान प्रत्यक्ष स्वैप तर्क}

वर्तमान प्रत्यक्ष पूल स्वैप:

\begin{enumerate}[leftmargin=*]
\item इनपुट और आउटपुट टोकन के लिए वैकल्पिक रजिस्टर की जांच करता है;
\item प्राप्त इनपुट को मापता है;
\item कॉन्फ़िगर किए गए कोटेटर से उद्धरण प्राप्त करता है या कच्चे इकाई समता लागू करता है;
\item वैकल्पिक लिमिटर के मुकाबले प्राप्त पूल टोकन बैलेंस की जांच करता है;
\item समूह शुल्क और किसी भी अतिरिक्त प्रोटोकॉल शुल्क की गणना करता है;
\item उपलब्ध आउटपुट इन्वेंट्री की जांच करता है;
\item प्रोटोकॉल शुल्क और आउटपुट को स्थानांतरित करता है और पूल शुल्क के लिए लेखांकन करता है; तथा
\item स्वैप घटनाओं का उत्सर्जन करता है।
\end{enumerate}
उद्धरण एक लेनदेन पैरामीटर है, न कि जारीकर्ता क्षमता का प्रमाण, चुकौती मूल्य, निष्पक्ष मूल्य, नकदी परिवर्तनीयता या गारंटी।

\subsection*{1.4 व्यापक उपयोग}
\addcontentsline{toc}{subsection}{1.4 व्यापक उपयोग}

प्रतिबद्धता पूल सामुदायिक आदान-प्रदान, उत्पादन, आपसी सहायता, सार्वजनिक कार्यक्रमों और अन्य जिम्मेदार संरचनाओं का समर्थन कर सकता है। एक अलग से प्रलेखित क्रेडिट उत्पाद गारंटी या चुकौती के साधन के रूप में वाउचर का उपयोग कर सकता था, लेकिन इसके लिए अतिरिक्त और लेनदेन विशिष्ट शर्तों की आवश्यकता होगी। एक सामान्य भेजने, पूल जमा, पूल स्वैप, भुनाने के लिए प्रस्तुति, वादे की पूर्ति या डिस्चार्ज स्वचालित रूप से ऋण या चुकाने नहीं है।

CPP के लिए जिम्मेदार एक्सचेंज और ऑडिटेबल लेनदेन रिकॉर्ड हैं, न कि अनुमानित चर्न। ऑन-चेन रिकॉर्ड अभी भी वास्तविक दुनिया में वादे की पूर्ति या सामाजिक प्रभाव साबित नहीं करते हैं।


\section*{2. लेखा परिवर्तनः परिसंपत्तियों से ट्रस्ट तक}
\addcontentsline{toc}{section}{2. लेखा परिवर्तनः परिसंपत्तियों से ट्रस्ट तक}

पारंपरिक वित्त \textbf{सम्पत्ति - देयता = इक्विटी}. CPP यह एक प्रतिबद्धता अर्थव्यवस्था के लिए रीफ्रेम करता हैः

\begin{itemize}[leftmargin=*]
\item \textbf{स्वीकृति क्षमताः} एक समूह या नेटवर्क अभी भी किसी जारीकर्ता की प्रतिबद्धताओं में से कितने को स्वीकार करने के लिए तैयार है।
\item \textbf{बकाया प्रतिबद्धताः} ऐसे वादे किए जाते हैं जो पूरे नहीं होते और दूसरों द्वारा रखे जाते हैं।
\item \textbf{वादे की पूर्ति क्षमताः} इन वादे को पूरा करने के लिए जारीकर्ता की सिद्ध क्षमता।
\end{itemize}
\subsection*{प्रतिबद्धता-क्षमता संबंध का उदाहरणः}
\addcontentsline{toc}{subsection}{प्रतिबद्धता-क्षमता संबंध का उदाहरणः}

स्वीकृति क्षमता - बकाया प्रतिबद्धताएं = शेष स्वीकृती क्षमता

यह अवधारणात्मक पहचान पूल जोखिम सीमाओं को सूचित कर सकती हैः असीमित जारी करने का मतलब असीम स्वीकृति नहीं है। यह एक बैलेंस पहचान या बयान नहीं है कि प्रत्येक वाउचर कानूनी रूप से ऋण या ऋण है।

\textbf{उदाहरण:} एक ट्रांसपोर्टर Vouchers में 100100 राइड्स'' जारी करता है। एक पूल एक बार में अधिकतम 40 राइड स्वीकार करता है. यदि 15 स्वीकृत राइड अप्रासंगिक रहते हैं, तो पूल को उस पॉलिसी के तहत 25 अतिरिक्त राइडों तक स्वीकार करने की क्षमता होती है। गणना स्वयं ऋण या गारंटी वादे की पूर्ति का निर्माण नहीं करती है।



\section*{3. वादे की पूर्ति, निकासी और विनिमय गतिविधि}
\addcontentsline{toc}{section}{3. वादे की पूर्ति, निकासी और विनिमय गतिविधि}

एक पूल स्वैप परिवर्तन जो होल्ड परिसंपत्तियों को संबोधित करता है। वाउचर वादे की पूर्ति तब होती है जब जारीकर्ता प्रकाशित प्रतिबद्धता को पूरा करता है. डिस्चार्ज पूर्ण इकाइयों को रिकॉर्ड करता है ताकि उन्हें फिर से प्रस्तुत नहीं किया जा सके। इन घटनाओं को ``सेटलमेंट के एक उपयोग में नहीं गिराया जाना चाहिए।''

प्रस्तावित माप ढांचे में निम्नलिखित के लिए अलग-अलग रिकॉर्ड्स का उपयोग किया जाता हैः

\begin{itemize}[leftmargin=*]
\item पूल स्वैप की पूर्ण मात्रा;
\item वैध प्रतिपूर्ति प्रस्तुतियाँ;
\item जारीकर्ता द्वारा पुष्टि की गई वादे की पूर्ति;
\item दायित्व की समाप्ति;
\item बकाया पात्र प्रतिबद्धताएं; तथा
\item मापा गया पूल इन्वेंट्री।
\end{itemize}
एक प्रस्तावित प्रतिबद्धता-निष्कासन गति एक अवधि के दौरान पूरा किए गए मूल्य की तुलना औसत बकाया पात्र प्रतिबद्धता मूल्य के साथ कर सकती है, लेकिन केवल तभी जब दोनों एक ही खुलासा मूल्य निर्धारण विधि का उपयोग करते हैं। टोकन आपूर्ति एक पर्याप्त संकेतक नहीं हैः इसमें जारीकर्ता सूची, समाप्त या अनुपलब्ध इकाइयां, परीक्षण और शेष शामिल हो सकते हैं जो बकाए तीसरे पक्ष की प्रतिबद्धता का प्रतिनिधित्व नहीं करते हैं.

एक अलग पूल स्वैप-कार्यात्मकता अनुपात मापी गई पूल इन्वेंट्री के साथ पूरा पूल स्वाप मूल्य की तुलना कर सकता है। यह विनिमय गतिविधि का वर्णन करता है, न कि जारीकर्ता की वादे की पूर्ति।

तरलता इन्वेंट्री की उपलब्धता में सुधार कर सकती है और आदान-प्रदान को आसान बना सकती है। यह गारंटी नहीं देती है कि जारीकर्ता प्रदर्शन करेंगे, कि योगदानकर्ता परिसंपत्तियों को पुनर्प्राप्त करेंगे, या कि एक पूल बोली भुनाने या नकदी मूल्य का प्रतिनिधित्व करती है। किसी भी तरलताओं के प्रदाता अधिकारों के लिए अलग प्रकाशित शर्तें आवश्यक हैं।

देखिये \href{https://docs.cosmolocal.credit/hi/white-paper/appendix-a-math-box}{परिशिष्ट ए} परिभाषाओं और \href{https://docs.cosmolocal.credit/hi/white-paper/appendix-c-kpi-definitions}{परिशिष्ट सी} प्रस्तावित माप विनिर्देश के लिए।


\section*{4. पुनर्नवीनीकरण योग्य वायदा शैली की प्रतिभूतियां}
\addcontentsline{toc}{section}{4. पुनर्नवीनीकरण योग्य वायदा शैली की प्रतिभूतियां}

एक अलग से प्रलेखित क्रेडिट सुविधा प्रतिभूतियों या चुकौती के साधन के रूप में फोंगेबल, हस्तांतरण योग्य वाउचर स्वीकार कर सकती है। यदि उसने ऐसा कियाः

\begin{itemize}[leftmargin=*]
\item पूल में प्रवेश से प्रतिभूतियों की अधिक खोज संभव हो सकती है;
\item अतिरिक्त वैध प्रस्तुति और वादे की पूर्ति विकल्प धनवापसी का समर्थन कर सकते हैं;
\item सीमाएं, इन्वेंट्री, मूल्यांकन, तरलता, जारीकर्ता प्रदर्शन और डिफॉल्ट जोखिम बरकरार रहेंगे; तथा
\item कोई मार्ग, वादे की पूर्ति, वापसी, मूल्य या वसूली की गारंटी नहीं दी जाएगी।
\end{itemize}
\subsection*{4.1 उत्पादक क्रेडिट लूप}
\addcontentsline{toc}{subsection}{4.1 उत्पादक क्रेडिट लूप}

एक भविष्य CLC- संगत सेवा केवल तभी उत्पादक ऋण का समर्थन कर सकती है जब पार्टियां एक अलग सुविधा स्थापित करें और लेनदेन को स्पष्ट रूप से वापस करने योग्य अग्रिम के रूप में प्रस्तुत करें। इसके लेनदेन की शर्तें क्रेडिटर और देनदार की पहचान करेगी और बताएगीः

\begin{itemize}[leftmargin=*]
\item अग्रिम और प्रतिपूर्ति की राशि;
\item परिपक्वता की तारीखें, ब्याज, शुल्क और अनुमत आश्वासन उपयोग;
\item असाइनमेंट और लेनदार या दावाधारक में कोई बदलाव;
\item भुगतान की लेखांकन के लिए वाउचर वादे की पूर्ति का मूल्यांकन और प्रमाण कैसे किया जाता है;
\item आंशिक धनवापसी, शेष राशि, default और discharge; तथा
\item लागू कानून के तहत उपलब्ध उपचार।
\end{itemize}
एक उदाहरणीय प्रवाह हैः

\begin{enumerate}[leftmargin=*]
\item उत्पादक या सेवा प्रदाता भविष्य में उत्पादन के लिए प्रतिबद्धता का प्रतिनिधित्व करने वाला एक वाउचर जारी करता है, जैसे दस टैक्सी सवारी, 50 किलोग्राम मक्का या दस श्रम घंटे।
\item एक पूल प्रबंधक उस वाउचर को स्वीकार करता है और पूल की विनिमय दर पद्धति, सीमाएं, शुल्क, इन्वेंट्री नियम और किसी भी स्पष्ट रूप से माना गया संरक्षण प्रकाशित करता है।
\item एक ऋणदाता या तरलता कार्यक्रम अलग से लिखित लेनदेन शर्तों के तहत अग्रिम प्रदान करता है और निर्माता के वाउचर को प्रतिभूति या सहमत वापसी साधन के रूप में स्वीकार करता है।
\item धारक वाउचरों को स्थानांतरित कर सकते हैं, उन्हें संगत पूल के माध्यम से आदान-प्रदान कर सकते है या उन्हें जारीकर्ता को प्रस्तुत कर सकता है।
\item जारीकर्ता द्वारा पुष्टि की गई वादे की पूर्ति वाउचर प्रतिबद्धता को पूरा करती है। यह केवल क्रेडिट शर्तों में निर्दिष्ट प्रमाण पर, मूल्य पर और साक्ष्य के आधार पर अलग-अलग ऋण शेष को कम करता है।
\item रिकॉर्ड वाउचर वादे की पूर्ति को क्रेडिट लेखा से अलग करते हैं और किसी भी आंशिक वापसी, शेष राशि, असाइनमेंट, डिफॉल्ट या डिस्चार्ज की पहचान करते हैं।
\end{enumerate}
इस ढांचे से वैचर और क्रेडिट दायित्वों के लिए अलग-अलग रिकॉर्ड बनाए रखते हुए सहमत वस्तुगत धनवापसी की अनुमति दी जा सकती है। यह एक साधारण हस्तांतरण, जमा, पूल स्वैप, चुकौती प्रस्तुति, वादे की पूर्ति या डिस्चार्ज से उत्पन्न नहीं होता है।

\textbf{उदाहरणीय उदाहरण:} एक उत्पादक को निर्दिष्ट मक्का वाउचरों के पुष्ट वादे की पूर्ति की पुष्टि करने की शर्तों के तहत 1,000 डॉलर का अग्रिम प्राप्त होता है, जिसमें कहा जाता है कि निर्दिष्ट मूल्य निर्धारण पद्धति का उपयोग करके धनवापसी राशि का श्रेय दिया जाता है। उधारकर्ता उन शर्तों द्वारा आवश्यक सबूत प्राप्त करने के बाद ही ऋण लागू करता है। यदि क्रेडिट शर्तें उस संबंध को नहीं करती हैं, तो वाउचर का अधिग्रहण, हस्तांतरण, आदान-प्रदान, प्रस्तुत करना या पूरा करना अग्रिम पर प्रभाव नहीं डालता है।


\section*{5. अलग-अलग पूल से एक संघीय नेटवर्क तक}
\addcontentsline{toc}{section}{5. अलग-अलग पूल से एक संघीय नेटवर्क तक}

वर्तमान Protocol v1.1.0 एक के माध्यम से प्रत्यक्ष निष्पादन का समर्थन करता है\texttt{SwapPool} और केवल उद्धरण प्रदान करता है\texttt{SwapRouter}. यह मल्टी-हॉप रूट, एचटीएलसी, एस्क्रो रूट या बैच नेटिंग या क्रॉस नेटवर्क क्लिअरिंग नहीं करता है।

इस अध्याय में प्रस्तावित किया गया है कि स्वतंत्र रूप से शासित पूल अपने स्वयं के प्रवेश, मूल्यांकन, सीमा, शुल्क, सूची, प्राधिकरण और शासन नियमों को छोड़ने के बिना कैसे समन्वय कर सकते हैं।

\subsection*{5.1 अलग-अलग विनिमय और वादे की पूर्ति उपाय}
\addcontentsline{toc}{subsection}{5.1 अलग-अलग विनिमय और वादे की पूर्ति उपाय}

फेडरेशन इन्वेंट्री तक पहुंच में सुधार कर सकता है, लेकिन यह वाउचर को मिलाकर जीवनचक्रों का आदान-प्रदान नहीं करेगा। कोई भी कार्यान्वयन इन घटनाओं को अलग से मापता हैः

\begin{enumerate}[leftmargin=*]
\item एक मार्ग उद्धृत किया गया है;
\item एक या अधिक पूल स्वैप श्रृंखला पर निष्पादित और निपटान करते हैं;
\item एक धारक जारीकर्ता को वाउचर इकाइयां प्रस्तुत करता है;
\item जारीकर्ता प्रतिबद्धता को पूरा करता है; तथा
\item पूरा किए गए इकाइयों को छोड़ दिया जाता है।
\end{enumerate}
अधिक उद्धृत या निष्पादित मार्ग अधिक वादे की पूर्ति साबित नहीं करते हैं। रिपोर्ट में कोहर्ट, अवधि, परिसंपत्तियां, मूल्यांकन विधि और समयशीर्षक, बहिष्करण, सुधार और श्रृंखला से बाहर साक्ष्य निर्दिष्ट किए जाएंगे।

\textbf{इलस्ट्रेटिव मार्ग:} एक स्कूल मक्का वाउचर रखता है लेकिन परिवहन वाउचर्स की आवश्यकता होती है। एक रूट सेवा संगत पूल इन्वेंट्री की पहचान करती है। निष्पादन केवल तभी सफल होगा यदि प्रत्येक अलग से अधिकृत हूप अपनी बोली सीमाओं, सीमाएं, शुल्क, सूची और नीति के भीतर बने रहे। परिणामस्वरूप स्वैप यह साबित नहीं करेंगे कि किसी भी जारीकर्ता ने बाद में अपने वाउचर प्रतिबद्धताओं को पूरा किया था।

\subsection*{5.2 प्रस्तावित रूटिंग और रिबालेंसिंग सेवाएं}
\addcontentsline{toc}{subsection}{5.2 प्रस्तावित रूटिंग और रिबालेंसिंग सेवाएं}

भविष्य की रूट सेवा दो अलग-अलग गतिविधियों का समर्थन कर सकती है।

\textbf{प्रतिभागी द्वारा शुरू की गई निष्पादन।} इनपुट और आउटपुट परिसंपत्तियों, एक राशि और उपयोगकर्ता प्रतिबंधों को देखते हुए, सेवा एक पथ की पहचान कर सकती है और निष्पादन की तैयारी करती है। प्रत्येक हॉप का अपना जिम्मेदार पूल, बोली, प्राधिकरण, शुल्क, सीमाएं, इन्वेंट्री और रसीद होगी। परमाणु बैच, एचटीएलसी और एस्क्रो भविष्य के निष्पादन विकल्प हैं, वर्तमान प्रोटोकॉल व्यवहार नहीं।

\textbf{ऑप्ट-इन पूल पुनः संतुलन.} पूल प्रबंधक इन्वेंट्री लक्ष्य, अनुमत प्रतिभागियों, परिसंपत्ति वर्गों, उद्धरण विचलन सीमाओं और प्रति अवधि सीमाओं को प्रकाशित कर सकता है। एक उत्तरदायी सेवा संगत चक्र या श्रृंखलाओं की खोज कर सकती है और केवल अधिकृत इरादों को निष्पादित करती है।

एक पूल प्रतिभागियों के मार्गों की अनुमति दे सकता है जबकि आउटबाउंड रिबैंसिंग से इनकार कर सकता है, या यह केवल चयनित परिसंपत्तियों, समकक्षों और राशि को सक्षम कर सकता था। प्रत्येक निष्पादित होप एक रसीद उत्पन्न करेगा, और किसी भी सेवा शुल्क का खुलासा पूल और प्रोटोकॉल शुल्क से अलग किया जाएगा।

\subsubsection*{5.2.1 संघ और सहकार्यशीलता}

स्वतंत्र तैनाती अपने स्वयं के रजिस्टर, इंटरफेस, रूट सेवाओं और नीति प्रोफाइल को संचालित कर सकती है जबकि संगत डेटा और प्राप्ति मानकों का चयन करती है। क्रॉस-प्रोफाइल निष्पादन तैनातन पर निर्भर रहेगा।

एक संगत प्रोफ़ाइलः

\begin{itemize}[leftmargin=*]
\item अपनी रजिस्ट्री जड़ों, सेवा प्रदाताओं, नियंत्रकों और लागू शर्तों की पहचान करें;
\item अनुमत और अस्वीकृत प्रतिभागियों, परिसंपत्तियों, एडाप्टरों और मार्गों का खुलासा करना;
\item प्रत्येक प्रतिभागी पूल के प्राधिकरण, सीमाओं, शुल्क और इन्वेंट्री प्रतिबंधों को लागू करें;
\item प्रति हॉप उद्धरण-से-प्राप्ति साक्ष्य बनाए रखें; तथा
\item शेष राशि या जारीकर्ता के दायित्वों को मिटाए बिना अन्यथा कार्यात्मक पूल को एक और रजिस्टर छोड़ने या चुनने की अनुमति दें।
\end{itemize}
संगतता उपलब्ध विनिमय मार्गों को बढ़ा सकती है और एक रजिस्ट्री या ऑपरेटर पर निर्भरता को कम कर सकती है।CLC App, GEF, या किसी अन्य पूल के लिए जिम्मेदार एक जारीकर्ता की वादे की पूर्ति.

\subsection*{5.3 प्रस्तावित नेटवर्क रैक और सेवा शुल्क मॉडल}
\addcontentsline{toc}{subsection}{5.3 प्रस्तावित नेटवर्क रैक और सेवा शुल्क मॉडल}

वर्तमान Protocol v1.1.0 एक पूल शुल्क और, जब कॉन्फ़िगर किया जाता है, तो एक प्रत्यक्ष पूल स्वैप पर एक अतिरिक्त प्रोटोकॉल शुल्क लेता है। ये मौजूदा शुल्क अलग रहते हैं।

भविष्य के कार्यक्रम को अलग से प्राप्त किया जा सकता हैः

\begin{enumerate}[leftmargin=*]
\item एक \textbf{नेटवर्क रैक}, भाग लेने वाले पूल द्वारा एकत्रित पूल शुल्क के एक घोषित हिस्से के रूप में परिभाषित किया गया है; तथा
\item एक \textbf{रूटिंग या सेवा शुल्क}, एक ज्ञात भविष्य की सेवा के लिए चार्ज किया गया।
\end{enumerate}
प्रस्तावित नेटवर्क रैक पूल शुल्क की गणना के बाद पूल की पूरी राशि पर लागू अतिरिक्त प्रतिशत नहीं है।\texttt{p}:

\[
\tau_p = f_p \cdot r_p
\]

जहां\texttt{f\_p} यह पूल शुल्क दर है और \texttt{r\_p} यह नेटवर्क कार्यक्रम के लिए आवंटित उस पूल शुल्क का प्रस्तावित हिस्सा है।

एक मापी अवधि के लिएः

\begin{itemize}[leftmargin=*]
\item \textbf{सकल पूल शुल्क} प्रत्येक पूल के वास्तविक संग्रह किए गए पूल शुल्क का योग है;
\item \textbf{नेटवर्क रैक की प्राप्ति} उन एकत्रित शुल्क के घोषित हिस्से हैं;
\item \textbf{सेवा शुल्क रसीदें} रूटिंग या सेवा शुल्क अलग से एकत्र किए जाते हैं; तथा
\item \textbf{कार्यक्रम शुल्क की प्राप्ति} समान नेटवर्क रैक प्राप्ति और सेवा शुल्क प्राप्ति।
\end{itemize}
वर्तमान प्रोटोकॉल शुल्क को तब तक शामिल नहीं किया जाता जब तक कि एक अलग रूप से अपनाई गई नीति वास्तविक प्रोटोकॉलर शुल्क प्राप्तियों को भविष्य के कार्यक्रम में कानूनी रूप से पुनर्निर्देशित न करे।

एक समग्र अनुमान के लिए, let:

\begin{itemize}[leftmargin=*]
\item \texttt{Q\_swap} परिभाषित समूह और अवधि के लिए निष्पादित पूल स्वैप का मूल्य होगा; तथा
\item \texttt{$\tau$} प्रस्तावित नेटवर्क रैक की प्रभावी दर और उस निष्पादित स्वैप मूल्य पर अलग से निर्धारित सेवा शुल्क होगी।
\end{itemize}
फिरः

\[
F \approx \tau \cdot Q_{\text{\latinfont swap}}
\]

यह एक विश्लेषणात्मक अनुमान है, राजस्व का वादा नहीं। प्रत्येक इनपुट के लिए एक निर्दिष्ट कोहॉर्ट, अवधि, इकाई, मूल्यांकन समय टिकट, बहिष्करण और सुधार नीति की आवश्यकता होती है।

\subsubsection*{5.3.1 नकद पात्र रसीदें और प्रकृतियां}

शुल्क नकदी के पात्र फंजिबल परिसंपत्तियों में या वाउचर और अन्य वस्तुनिष्ठ परिसंपत्ति में आ सकते हैं। वस्तुनिहित रसीदें नकद व्यय या कवरेज दावों का स्वचालित रूप से भुगतान नहीं कर सकती हैं। किसी भी विनिमय या रूपांतरण के लिए प्राधिकरण, उपलब्ध इन्वेंट्री, खुलासा किए गए स्थानों, सीमाओं और वास्तविक निष्पादन की आवश्यकता होगी।

\texttt{$\chi$} को शुल्क प्राप्तियों का वास्तविक हिस्सा माना जाए जो पॉलिसी प्रतिबंधों, असफल रूपांतरणों और स्लिपिंग के बाद नकद पात्र है। नकद-उपयोग योग्य प्राप्तियां हैंः

\[
F_{\text{\latinfont cash}} \approx \chi \cdot F
\]

बजट और ब्रेक-ईवेंस विश्लेषण में प्राप्त \texttt{F\_cash} का उपयोग किया जाएगा, न कि सकल कोटेड शुल्क या इन-चीज इन्वेंट्री के नाममात्र मूल्य का। भविष्य के कार्यक्रम में कुल पूल शुल्क, नेटवर्क रेक रसीद, सेवा शुल्क, परिसंपत्ति संरचना, रूपांतरण परिणाम और नकदी-उपयोगी रसीदों की अलग से रिपोर्ट की जाएगी।

\subsection*{5.4 प्रस्तावित तरलता कार्यक्रम}
\addcontentsline{toc}{subsection}{5.4 प्रस्तावित तरलता कार्यक्रम}

भविष्य में, एक अलग से प्रलेखित तरलता कार्यक्रम निर्दिष्ट पूल या रूटिंग सेवाओं के लिए संपत्ति आवंटित कर सकता है। वर्तमान \texttt{SwapPool} अनुबंध पूल शेयरों को न तो मुद्रांकन करते हैं और न ही भुगतान, निकासी, इनाम, शासन या लाभ अधिकार स्वचालित रूप से बनाते हैं।

कोई भी कार्यक्रम प्रकाशित करेगाः

\begin{itemize}[leftmargin=*]
\item उत्तरदायी इकाई और भाग लेने वाली पूल प्रबंधक;
\item योगदान की गई संपत्तियां और क्या हस्तांतरण वापस करने योग्य, निकासी योग्य, दान या उपहार है;
\item रखरखाव और तकनीकी-नियंत्रण व्यवस्थाएं;
\item अनुमत उपयोग, सीमाएं, लॉकअप, वापसी गेट और हानि आवंटन;
\item शुल्क या प्रोत्साहन पात्रता और क्या राशि शून्य हो सकती है;
\item रिपोर्टिंग, संघर्ष, शिकायत और उपचार; तथा
\item शेष परिसंपत्तियों और दायित्वों का प्रवासन, समाप्ति और उपचार।
\end{itemize}
भौतिक जोखिमों में इन्वेंट्री शामिल है जिसे बदलना या पूरा करना मुश्किल है, कम नकदी पात्रता, जारीकर्ता का गैर-प्रदर्शन, अनुबंध या प्रदाता विफलता, शासन परिवर्तन और बाहर निकलने पर प्रतिबंध शामिल हैं। सीमाएं, भंडार, रसीदें और डैशबोर्ड कुछ जोखिमों को कम कर सकते हैं या प्रकट कर सकते हैंः वे नुकसान को समाप्त नहीं करते हैं।

एक्स-पोस्ट एनालिटिकल मीट्रिक के लिएः

\begin{itemize}[leftmargin=*]
\item \texttt{$\phi$} कार्यक्रम की स्वीकृत शर्तों के तहत आवंटित कार्यक्रम शुल्क प्राप्तियों का वास्तविक अंश होगा; तथा
\item \texttt{K} एक निर्दिष्ट विधि के तहत कार्यक्रम द्वारा कवर किए गए परिसंपत्तियों का मापा मूल्य होगा।
\end{itemize}
फिरः

\[
\text{\latinfont FeeFlow}_{LP} \approx \frac{\phi \cdot F}{K} = \frac{\phi \cdot \tau \cdot Q_{\text{\latinfont swap}}}{K}
\]

यह मीट्रिक प्रति मापा गया कार्यक्रम परिसंपत्ति के लिए प्राप्त शुल्क प्रवाह का वर्णन करता है। यह APY, एक पूर्वानुमान, एक लाभांश या गारंटीकृत रिटर्न नहीं है। रिपोर्टों में स्वैप वॉल्यूम, भुनाने के लिए प्रस्तुति, जारीकर्ता वादे की पूर्ति, डिस्चार्ज, होल्डिंग अवधि, नुकसान, निकासी और शुल्क रसीद को अलग रखा जाएगा।


\section*{6. प्रस्तावित नेटवर्क स्तर की तरलता और शासन}
\addcontentsline{toc}{section}{6. प्रस्तावित नेटवर्क स्तर की तरलता और शासन}

ऐतिहासिक Sarafu Network के अनुभव ने दो व्यापक डिजाइन प्रश्नों पर शोध को प्रेरित कियाः

\begin{enumerate}[leftmargin=*]
\item स्वतंत्र रूप से शासित पूल लिक्विडिटी और रूटिंग सेवाओं का समन्वय कैसे कर सकते हैं; तथा
\item कैसे साझा सेवाएं बढ़ती भागीदारी के साथ जवाबदेही बनाए रख सकती हैं।
\end{enumerate}
एक प्रस्तावित CLC डिजाइन \textbf{प्रस्तावित CLC नेटवर्क पूल } नेटवर्क स्तर पर क्लियरिंग और बजट समन्वय व्यवस्था के रूप में \textbf{ प्रस्तावित CLC शासन टोकन}. वर्तमान ऐप या Protocol v1.1.0 में कोई भी नहीं है. अन्य तैनाती सहकारी संस्थाओं, सार्वजनिक एजेंसियों, संघों, बहुसंस्थाओं, सेवा ऑपरेटरों या किसी भी प्रस्ताव को स्वीकार किए बिना अन्य जिम्मेदार संरचनाओं का उपयोग कर सकती हैं।

\subsection*{6.1 डाउनसाइड कंट्रोल और वैकल्पिकता}
\addcontentsline{toc}{subsection}{6.1 डाउनसाइड कंट्रोल और वैकल्पिकता}

वर्तमान Protocol v1.1.0 पूल टोकन बैलेंस कैप और इन्वेंट्री चेक लागू कर सकता है। ये नियंत्रण चयनित अनुबंध कार्यों को प्रतिबंधित करते हैं; वे जारीकर्ता के गैर-प्रदर्शन, मूल्य हानि, कुंजी हानि, अनुबंध विफलता या किसी अन्य नुकसान को नहीं रोकते हैं।

भविष्य की तैनाती में सर्किट ब्रेकर, टाइमलॉक, रिजर्व, गारंटी या बीमा शामिल हो सकते हैं। किसी भी सुरक्षा के लिए एक पहचाने गए जिम्मेदार पक्ष, कवर किए गए घटना, वित्त पोषण, सीमा, बहिष्करण, अवधि, सबूत, दावा प्रक्रिया और लागू शर्तों की आवश्यकता होगी।

प्रस्तावित शासन को रजिस्ट्री और सेवा शक्तियों को संकीर्ण और समीक्षा योग्य बनाए रखना चाहिए। सामान्य कार्यों में सूचना, देरी, कारण प्रकाशन, अपील और सुधार प्रक्रियाओं का उपयोग किया जा सकता है। आपातकालीन कार्यों को घटना रिकॉर्ड तैयार करना चाहिए और समाप्त हो जाना चाहिए या अपनाई गई नीति के तहत समीक्षा प्राप्त करनी चाहिए।

संगत रूट डिस्कवरी पूल में अधिक एक्सचेंजों को संभव बना सकती है। मल्टी-हॉप निष्पादन और नेटिंग के लिए अतिरिक्त सॉफ़्टवेयर, प्राधिकरण, सीमाओं, विफलता प्रबंधन और शासन की आवश्यकता होगी। अधिक उद्धृत या निष्पादित स्वैप अपने आप में वाउचर वादे की पूर्ति या सामाजिक प्रभाव साबित नहीं करेंगे।


\section*{7. प्रस्तावित CLC प्रबंधन और शासन संपत्ति}
\addcontentsline{toc}{section}{7. प्रस्तावित CLC प्रबंधन और शासन संपत्ति}

इस अध्याय में भविष्य के शासन और तरलता समन्वय का एक वैकल्पिक डिजाइन प्रस्तुत किया गया है।CLCशासन टोकन,stCLC, sCLC, शासन वॉल्ट, बीमा फंड, जलप्रपात, गेज, जनादेश, शुल्क-क्रेडिट विंडो और बाहरी बाजार कार्यक्रम वर्तमान में भाग नहीं लेते हैंCLC App याProtocol v1.1.0. प्रत्येक को कार्यान्वयन, उत्तरदायी ऑपरेटर, प्रकाशित नीति और शर्तें, सुरक्षा समीक्षा और लागू कानूनी अनुमोदन की आवश्यकता होगी।

CPP को इस टोकन मॉडल की आवश्यकता नहीं है। संगत तैनाती इसके बजाय सहकारी, सार्वजनिक एजेंसियों, संघों, मल्टी-सिग्स, सेवाओं या अन्य जिम्मेदार संरचनाओं का उपयोग कर सकती है।

\subsection*{7.1 उद्देश्य}
\addcontentsline{toc}{subsection}{7.1 उद्देश्य}

यदि इसे अपनाया जाता है, तो प्रस्तावित शासन स्तरः

\begin{itemize}[leftmargin=*]
\item साझा रजिस्टरों और मार्ग सेवाओं का समन्वय करना;
\item स्वतंत्र रूप से विनियमित पूल में स्वीकृत तरलता जनादेशों का आवंटन;
\item एक वैकल्पिक, स्पष्ट रूप से लक्षित, वित्त पोषित कवरेज कार्यक्रम का प्रबंधन करना;
\item सामान्य बुनियादी ढांचे और अवलोकनशीलता को बनाए रखना;
\item प्रकाशित पात्रता और गेमिंग विरोधी नियमों के तहत योगदानकर्ताओं को पहचानें; तथा
\item भागीदारी बढ़ने के साथ समीक्षा, अपील, पारदर्शिता और विश्वसनीय बाहर निकलने को बनाए रखना।
\end{itemize}
\subsection*{7.2 प्रस्तावित शासन-सम्पत्ति मॉडल}
\addcontentsline{toc}{subsection}{7.2 प्रस्तावित शासन-सम्पत्ति मॉडल}

डिजाइन में तीन अलग-अलग प्रस्तावित संपत्ति शामिल हैंः

\begin{enumerate}[leftmargin=*]
\item \textbf{प्रस्तावित CLC शासन टोकनः} स्वीकृत जारी, रखरखाव, मतदान और अनुपालन नियमों के तहत एक हस्तांतरित आधार शासन परिसंपत्ति।
\item \textbf{stCLC:} प्रस्तावित CLC शासन टोकन के लॉक होने पर मुद्रित एक गैर-परिवर्तनीय वोट-एस्क्रो रसीद। यह समय सीमाबद्ध शासन शक्ति का प्रतिनिधित्व करेगा और खुलासा किए गए प्रतिनिधिमंडल की अनुमति दे सकता है।
\item \textbf{sCLC:} पॉलिसी के तहत जारी किए गए एक युग-व्यापक प्राधिकरण या प्रोत्साहन परिसंपत्ति। यह सीमाबद्ध शुल्क-क्रेडिट पहुंच या अनुमोदित तरलता और रूटिंग प्रोत्साहन का समर्थन कर सकता है और अवधि के अंत में समाप्त हो सकता है या जलाया जा सकता है।
\end{enumerate}
न हीstCLCऔर न हीsCLCएक नीति अपनाई गई है, जो कि एक इक्विटी, लाभांश, शेष राशि का दावा, वापसी का वादा या सामुदायिक वाउचर हो सकता है।sCLCकिसी भी युग में निष्पादन और शुल्क-क्रेडिट तक पहुंच शून्य है।

उपयोग से पहले शासन लॉकअप, न्यूनतम अवधि, शीतलन, प्रतिनियुक्ति, एकाग्रता सीमाएं और महत्वपूर्ण कार्रवाई की सीमाएँ प्रकाशित की जाएंगी। अनलॉक किए गए प्रस्तावित CLC शासन टोकन इस डिजाइन के तहत मतदान नहीं करेंगे।

\subsubsection*{7.2.1 आपूर्ति और आवंटन का चित्रण}

CLC गवर्नेंस टोकन की प्रस्तावित इलस्ट्रेटिव कुल आपूर्ति \textbf{500,000,000}. यह एक डिजाइन पैरामीटर है, न कि चालू जारी, प्रस्ताव, मूल्यांकन या तरलता का वादा।

एक स्पष्ट आवंटन हैः

\begin{itemize}[leftmargin=*]
\item \textbf{15\% --- Grassroots Economics Foundation:} स्थायी रूप से सट्टेबाजी, गैर हस्तांतरण योग्य और प्रकाशित हस्ताक्षरकर्ता और रोटेशन नियमों के तहत एक पहचान GEF मल्टीसिग से बंधा।
\item \textbf{15\% --- कोर टीम और शुरुआती साझेदारः} मतदान और हस्तांतरण के नियमों का पहले से खुलासा करते हुए, 24 महीने की नीति-निर्धारित चट्टान और रैखिक अधिग्रहण अवधि के दौरान दांव लगाया गया।
\item \textbf{30\% --- निजी अनुदान:} अनुमोदित नेटवर्क तरलता प्रदान करने वाले शुरुआती योगदानकर्ताओं के लिए मान्यता, 24 महीने की नीतिगत अधिग्रहण के साथ।
\item \textbf{40\% --- सार्वजनिक तरलता भंडारः} वैकल्पिक बाह्य बाजार और सुलभता कार्यक्रमों के लिए समय-बंद शासन वॉल्ट में आयोजित किया जाता है।
\end{itemize}
स्पष्ट सार्वजनिक तरलता नियंत्रण के अंतर्गतः

\begin{itemize}[leftmargin=*]
\item बाहरी स्थानों पर एक ही समय में कुल आपूर्ति का 10\% से अधिक सक्रिय नहीं होगा;
\item प्रत्येक तैनाती 90 दिनों के बाद समाप्त हो जाती है जब तक कि इसे अपनाई गई प्रक्रिया के माध्यम से नवीनीकृत नहीं किया जाता;
\item लिक्विडिटी पोजीशन और कोई भी पोजिशन टोकन खुला शासन नियंत्रण के अधीन रहेंगे; तथा
\item प्रस्तावित CLC शासन टोकन लिक्विडिटी पदों में उपयोग किए जाने वाले अन्य मतदान टोकन के समान नियमों के तहत लॉक किए बिना वोट नहीं करेंगे।
\end{itemize}
भविष्य में लॉन्च एक खुलासा किए गए मल्टीसिग और संक्रमण के साथ केवल स्वतंत्र रूप से अनुमोदित कठोरता मील के पत्थरों, जैसे कि स्वतंत्र लेखा परीक्षाओं, निगरानी, घटना रनबुक, और परीक्षण विराम और फोर्क प्रक्रियाओं के बाद शुरू हो सकता है। प्रत्येक अपनाया गया पैरामीटर और परिवर्तन प्रक्रिया को अलग से प्रकाशित किया जाएगा।

\subsubsection*{7.2.2 उपलब्धता और दान के चरण}

भविष्य के अनुदान कार्यक्रम में चरणबद्ध योगदान खिड़कियों का उपयोग किया जा सकता है और आवंटन और बजट के लिए एक प्रकाशित संदर्भ मूल्य। यह संदर्भ बाजार मूल्य, मूल्यांकन, वापसी या बाहर निकलने का वादा नहीं होगा।

योगदान की शर्तें बताती हैं कि क्या परिसंपत्तियां दान, दान, निकासी योग्य या वापस करने योग्य हैं; उन्हें कौन नियंत्रित करता है; उनका अनुमत उपयोग; लॉकअप; हानि आवंटन; रिपोर्टिंग; और बाहर निकलने की स्थिति। बड़े योगदानों का सामना मतदान सीमाओं या शासन वजन में कमी से हो सकता है।

किसी भी बाहरी बाजार तरलता के लिए स्थानों, जिम्मेदार ऑपरेटरों, परिसंपत्तियों, सीमाओं, समय निर्धारण, संघर्षों, जोखिम नियंत्रण, रिपोर्टिंग और समाप्ति की पहचान करने के लिए एक अलग निर्णय की आवश्यकता होगी। नीति मूल्य को लक्षित या गारंटी नहीं देगी।

\subsubsection*{7.2.3 प्रस्तावित प्रभाव रोपण कार्यक्रम}

भविष्य के कार्यक्रम में गवर्नेंस वॉल्ट का एक हिस्सा उन योगदानकर्ताओं को आवंटित किया जा सकता है जो निर्दिष्ट प्रतिबद्धता पूल में संपत्ति डालते हैं और एक्सचेंज तक पहुंच और प्रमाणित जारीकर्ता वादे की पूर्ति में मापनीय सुधार करते हैं।

एक स्वीकार की जाने वाली पात्रता नीति:

\begin{enumerate}[leftmargin=*]
\item अनुमोदित पूल, संपत्ति, न्यूनतम अवधि और लॉकअप की शर्तों की पहचान करें;
\item निष्पादित स्वैप साक्ष्य और अलग से प्रमाणित प्रस्तुति, वादे की पूर्ति और सप्लाई का उपयोग करके उत्पादक योगदान को परिभाषित करें;
\item अकेले लॉक किए गए कुल मूल्य के बजाय मार्जिनल योगदान को मापें;
\item स्व-विक्री और परिपत्र धोने की गतिविधि को बाधित करना;
\item प्रति इकाई सीमाएं और घटती रिटर्न लागू करें; तथा
\item अंतिम आवंटन से पहले एक अवलोकन, विवाद और सुधार विंडो प्रदान करें।
\end{enumerate}
कोई भी अनुदान कार्यक्रम के प्रकाशित शर्तों, अधिग्रहण, पात्रता, अनुपालन और हानि नियमों के अधीन रहेगा।

\subsection*{7.3 प्रस्तावित मतदान, प्रोत्साहन और शुल्क तक पहुंच}
\addcontentsline{toc}{subsection}{7.3 प्रस्तावित मतदान, प्रोत्साहन और शुल्क तक पहुंच}

इस डिजाइन के तहत,stCLCएक क्युरेटरी गेज सिस्टम वोटों को सीमित संख्या में अनुवाद कर सकता है।sCLCपात्र तरलता या रूटिंग ऑपरेटरों के लिए प्रोत्साहन।

मूल्यांकन नीति निष्पादित स्वैप और जारीकर्ता वादे की पूर्ति को अलग रखती है। यह उद्धृत नहीं बल्कि निष्पाद किए गए मार्गों, अनसुलझे प्रस्तुतियों, स्वयं-विनिमय या उपयोगी गतिविधि के प्रमाण के बिना कुल मूल्य को लॉक करेगी।

एक स्वीकृत शासन प्रक्रिया में एक युग शुल्क-क्रेडिट बजट भी प्रकाशित किया जा सकता है।stCLCधारकों को प्राप्त हो सकता हैsCLCसीमाबद्ध, समय-सीमित स्वैप को अनुमति देना नामित शुल्क धारक पूल से अनुदानित परिसंपत्तियों या कार्यक्रमों में। यह उपलब्ध सूची तक नीतिगत अभिगम होगा, निष्क्रिय आय या शुल्क धारन पूल का स्वामित्व नहीं।

\subsubsection*{7.3.1 सेवा शुल्क लागू करना और प्रोफ़ाइल चुनना}

भविष्य के रजिस्ट्री प्रोफ़ाइल में भाग लेने वाले पूल या रूट सेवाओं को एक संगत शुल्क एडाप्टर या हुक का उपयोग करने की आवश्यकता हो सकती है। यदि आवश्यक रसीद शुल्क संग्रह का प्रमाण नहीं देती है तो प्रोफाइल खोज, रूटिंग या कार्यक्रम समर्थन से इनकार कर सकता है।

जैसे नाम \textbf{PoolFactory } या \textbf{ FeeHook} संभावित भविष्य के मॉड्यूलों का वर्णन करें; वे Protocol v1.1.0 अनुबंध नहीं हैं। किसी भी कार्यान्वयन में कोड, पते, नियंत्रक, प्राप्तकर्ता, दरें, पात्र लेनदेन, विफलता प्रबंधन और लेखा परीक्षा की स्थिति का खुलासा किया जाएगा।

प्रवर्तन प्रोफाइल-विशिष्ट होगा। एक प्रोफ़ाइल द्वारा बाहर रखा गया पूल स्थानीय रूप से काम करना जारी रख सकता है या यदि उसके अनुबंध और निर्भरताएं कार्यात्मक रहें तो अन्य संगत रजिस्टर में शामिल हो सकता है। ग्राहकों को उपलब्ध प्रोफाईलों का खुलासा करना चाहिए और प्राधिकरण से पहले लागू शुल्क दिखाना चाहिए।

\subsubsection*{7.3.2 sCLC बजट और वैकल्पिक बाहरी फ्लोटिंग नियंत्रण}

उच्च प्राथमिकता वाले बजटों को अपनाए जाने के बाद, भविष्य की प्रक्रिया में एक युग शुल्क-क्रेडिट बजट प्रकाशित किया जा सकता है\texttt{F\_epoch} एक पॉलिसी प्रतिभागी सीमा के अनुपात में आवंटित कर सकती हैstCLCमतदान शक्तिः

\[
\text{\latinfont limit}_{\text{\latinfont user,epoch}} = F_{\text{\latinfont epoch}} \times \frac{\text{\latinfont stCLC}_{\text{\latinfont user}}}{\text{\latinfont stCLC}_{\text{\latinfont total}}}
\]

प्रत्येक खिड़की पर अनुमति सूची, कैप, इन्वेंट्री, पात्रता, घटना नियम और लागू कानून के अधीन रहेगा।

एक अलग नीति प्रस्तावित CLC शासन टोकन के सीमित, समय-वजन प्राप्त करने की अनुमति दे सकती है केवल बाहरी शासन-हमलों की मापी सतह को कम करने के लिए। यह आवश्यक होगाः

\begin{itemize}[leftmargin=*]
\item एक निर्दिष्ट अवधि के दौरान बाहरी फ्लोट पर आधारित प्रकाशित ट्रिगर;
\item प्राप्त, पात्र शुल्क और आयोजन स्थल की मात्रा का अधिकतम हिस्सा;
\item सटीक समय घोषणा के बिना समय-वजनित निष्पादन;
\item यदि अपनाए गए कवरेज या परिचालन सीमाओं को पूरा नहीं किया जाता है तो स्वचालित रोकना;
\item प्राप्त टोकन को एक खुलासा किए गए गैर-मतदाता खाते में सेवानिवृत्त करना या रखना; तथा
\item एक स्पष्ट कथन कि कार्यक्रम किसी मूल्य को लक्षित या गारंटी नहीं देता है और यह जारीकर्ता के वादे की पूर्ति को प्रभावित नहीं करता है।
\end{itemize}
पॉलिसी अनिश्चित काल के लिए अक्षम रह सकती है।

\subsubsection*{7.3.3 पोर्टफोलियो पूल और प्रमाणपत्र}

ए \textbf{पोर्टफोलियो पूल} एक साधारणप्रतिबद्धता पूलएक मिशन या सेवा क्षेत्र के आसपास संचालित किया गया है।पूल प्रबंधकयह एक व्यक्ति, सहकारी, सामुदायिक समूह, सार्वजनिक एजेंसी, संघ, मल्टीसिग, सेवा ऑपरेटर या अन्य जिम्मेदार संरचना हो सकती है।

सहायता निम्नलिखित के माध्यम से प्राप्त की जा सकती हैः

\begin{enumerate}[leftmargin=*]
\item पूल की प्रकाशित शर्तों के तहत प्रत्यक्ष योगदान;
\item स्वीकृत शासन प्रक्रिया के माध्यम से अनुमोदित समय-सीमित तरलता जनादेश; या
\item sCLC-निर्देशित अभिगम जब उस तंत्र को सक्षम किया गया हो।
\end{enumerate}
पोर्टफोलियो पूल्स स्वतंत्र रूप से शासित रहेंगे और साझा या स्वतंत्र रजिस्टरों का उपयोग कर सकते हैं।

प्रमाणपत्र प्रवेश या जोखिम उपचार को सूचित कर सकते हैं, लेकिन शुल्क, लाभ या अवशिष्ट परिसंपत्तियों के लिए अधिकार पैदा नहीं करेंगे। एक तैनाती का उपयोग किया जा सकता हैः

\begin{itemize}[leftmargin=*]
\item \textbf{रजिस्टर-बंद प्रमाणपत्र} एक प्रकाशित विधि का उपयोग करने वाले एक पहचाने गए सत्यापितकर्ता से; या
\item \textbf{सत्यापन सेवा वाउचर} एक घोषित लेखा परीक्षा या मूल्यांकन करने के लिए प्रतिपूर्ति योग्य प्रतिबद्धता का प्रतिनिधित्व करना।
\end{itemize}
पूल यह खुलासा करेगा कि एक प्रमाणपत्र प्रवेश, विनिमय दर के तरीकों, सीमाओं या पात्रता को कैसे प्रभावित करता है और त्रुटियों, समाप्ति, संघर्षों और अपीलों को कैसे संभाला जाता है।

\subsection*{7.4 प्रस्तावित झरना और बजट}
\addcontentsline{toc}{subsection}{7.4 प्रस्तावित झरना और बजट}

प्रस्तावित जलप्रपात में केवल प्राप्त किए गए, एक स्वीकृत शासन प्रक्रिया के तहत पात्र प्राप्तियों को आवंटित किया जाएगा।

\begin{itemize}[leftmargin=*]
\item \textbf{नकद पात्र संपत्ति (\texttt{E\_cash}):} सूचीबद्ध फ्यूजबल परिसंपत्तियां जिन्हें नकदी में मुद्रीकृत दायित्वों के लिए पॉलिसी के तहत इस्तेमाल या परिवर्तित किया जा सकता है; तथा
\item \textbf{प्रकृति में सम्पत्ति (\texttt{E\_kind}):} वाउचर या अन्य संपत्तियां जो नेटवर्क में विनिमय का समर्थन कर सकती हैं, लेकिन जब तक वास्तव में परिवर्तित नहीं होती हैं तब तक नकद मुद्रा कवर या परिचालन दायित्वों के लिए गिनती नहीं करती हैं।
\end{itemize}
एक रूपांतरण नीति में अधिकृत ऑपरेटरों, परिसंपत्तियों, स्थानों, मूल्य स्रोतों, समय खिड़कियों, स्लाइडिंग सीमाओं, कैप, संघर्षों, रिकॉर्ड और घटना प्रक्रियाओं की पहचान होगी। ट्रेजरी रूपांतरण किसी जारीकर्ता के वाउचर दायित्व या पूल की विनिमय दर पद्धति को निर्धारित नहीं करेगा।

एक स्पष्ट प्राथमिकता क्रम हैः

\begin{enumerate}[leftmargin=*]
\item \textbf{वित्त पोषित कवरेज लक्ष्य:} परिभाषित कवर की गई घटनाओं के लिए अपनाए गए रिजर्व या प्रस्तावित बीमा निधि को अपने खुलासा किए गए लक्ष्य के अनुसार बनाएँ।
\item \textbf{कोर ऑपरेशनः} कानूनी कार्य, शिक्षा, संचार, बुनियादी ढांचा, लेखा परीक्षा और अवलोकन के लिए एक सीमित, खुलासा किया गया बजट निधि।
\item \textbf{लिक्विडिटी जनादेशः} सीमाओं और सूर्यास्त की समीक्षा के तहत घोषित पूल या मार्ग सेवा उद्देश्यों के लिए अनुमोदित परिसंपत्तियों को आवंटित करें।
\item \textbf{वैकल्पिक बाहरी फ्लोटिंग नियंत्रणः} केवल तभी धनराशि जमा की जाती है जब उसके प्रकाशित ट्रिगर और उच्च प्राथमिकता वाले सीमाओं को पूरा किया जाता है; अन्यथा शून्य आवंटित करें।
\item \textbf{प्रस्तावित शुल्क-ऋण बजटः} किसी भी अनुमोदित शेष राशि को निर्दिष्ट शुल्क धारक पूल में आवंटित करें और \texttt{F\_epoch} प्रकाशित करें, जो शून्य हो सकता है।
\end{enumerate}
बजट निर्णयों में पुष्ट वादे की पूर्ति, कवर किए गए जोखिम, पात्र भंडार, सीमा उपयोग, निष्पादित मार्ग, एकाग्रता, घटनाओं और गारंटर के प्रदर्शन पर विचार किया जा सकता है। प्रत्येक मीट्रिक परिशिष्ट सी में साक्ष्य और कोहर्ट नियमों का पालन करेगा।

एक संभावित योगदानकर्ता प्रवाह होगाः

\begin{enumerate}[leftmargin=*]
\item योगदानकर्ता अलग से अपनाए गए अनुदान या कार्यक्रम शर्तों के तहत पात्र परिसंपत्तियों को स्थानांतरित करते हैं;
\item पात्र प्रतिभागियों को प्रस्तावित लॉक प्राप्त होता है और यदि अनुमति दी जाती हैCLCप्राप्त करने के लिए शासन टोकनstCLC;
\item पानी के झरने में प्राप्त किए गए पात्र रैक या सेवा शुल्क की प्राप्तियां प्रवेश करती हैं;
\item जलप्रपात निधियों ने क्रमबद्ध प्राथमिकताएं अपनाई हैं; तथा
\item केवल एक सक्षम वॉटरफ़ॉल के बाद नीति sCLC जारी करती है या एक सीमा शुल्क-क्रेडिट विंडो खोलती है।
\end{enumerate}
किसी भी कदम से स्वामित्व, वापसी, निकासी, कवरेज, इनाम या शासन के अधिकारों का निर्माण नहीं होता है जो लागू शर्तों में स्पष्ट रूप से बताए गए हैं।


\section*{8. तकनीकी दायरा और विकास}
\addcontentsline{toc}{section}{8. तकनीकी दायरा और विकास}

इस अध्याय में वैकल्पिक या प्रस्तावित कार्य का वर्णन किया गया है। यह वर्तमान में मौजूद होने की गारंटी वाली विशेषताओं की सूची नहीं हैCLC App याProtocol v1.1.0.

संभावित कार्यक्षेत्रों में निम्नलिखित शामिल हैंः

\begin{itemize}[leftmargin=*]
\item रजिस्टर, बोली, सीमा, शुल्क और इन्वेंट्री डिस्कवरी के साथ संगत पूल पर निष्पादन रूटिंग;
\item जब परमाणु निपटान अनुपलब्ध हो, तो समय पर लॉक किए गए एस्क्रो या HTLC एडाप्टर क्रॉस-डोमेन निष्पादन के लिए;
\item छोटे या व्यक्तिगत पूल के लिए इंटरफेस और नीतिगत उपकरण;
\item वाउचर, पूल, विनिमय दर पद्धति, सीमाओं, नियंत्रकों और शुल्क के लिए लेखापरीक्षण योग्य रजिस्टर;
\item तैनाती-विशिष्ट भुगतान प्रदाता कनेक्टर्स, चेकआउट प्रवाह और पात्रता नियंत्रण;
\item पुनः संतुलन और भुगतान लिक्विडिटी के लिए बाहरी तरलता स्थलों से फंगल एसेट कनेक्शन; तथा
\item नीतिगत कवरेज, परिचालन लागत या तरलता जनादेशों के लिए राजकोषीय रूपांतरण।
\end{itemize}
बाह्य बाजार की कीमतें निर्धारित नहीं करतीं कि एक जारीकर्ता वाउचर शर्तों के तहत कितना कर्जदार है। एक पूल फंजबल परिसंपत्ति के लिए एक संरक्षित बाहरी संदर्भ का उपयोग कर सकता है, लेकिन इसकी प्रकाशित विनिमय दर पद्धति, सीमाएं, शुल्क और इन्वेंट्री उसके उद्धरणों को नियंत्रित करेगी।

\subsection*{8.1 प्रस्तावित मार्ग-सेवा और SDK मानदंड}
\addcontentsline{toc}{subsection}{8.1 प्रस्तावित मार्ग-सेवा और SDK मानदंड}

\textbf{खोज.} प्रस्तावित रूट सेवा में परिसंपत्तियों के प्रवेश, विनिमय दर के तरीकों, सीमाओं, शुल्क, इन्वेंट्री, घटनाओं और नियंत्रक की जानकारी के लिए पहचाने गए रजिस्टरों से पूछताछ होगी। कैश किए गए रिकॉर्ड में ताज़ापन सीमाएं और स्रोत पहचानकर्ता शामिल होंगे।

\textbf{नेटवर्क प्रोफाइल।} एक ग्राहक एक से अधिक रजिस्ट्री रूट या नीति प्रोफाइल का समर्थन कर सकता है। यह प्रतिभागी को बताएगा कि कौन सा प्रोफ़ाइल, समकक्ष, एडाप्टर, प्रतिबंध और जिम्मेदार सेवा ऑपरेटर उद्धरण का उपयोग करता है। एक क्रॉस-प्रोफाइल मार्ग को प्रत्येक लागू हॉप की शर्तों को पूरा करने की आवश्यकता होगी।

\textbf{पथ नीति।} एक जिम्मेदार ऑपरेटर असुरक्षित निर्भरता या प्रतिभागियों को बाहर कर सकता है और मार्ग स्तर की सीमाओं, ताजापन आवश्यकताओं और स्वास्थ्य मानदंडों को लागू कर सकता हैः ये संकेत निर्णय का समर्थन करेंगे; वे वादे की पूर्ति या हानि से सुरक्षा की गारंटी नहीं देंगे।

\textbf{शुल्क और सीमाएं।} एक उद्धरण में पूल शुल्क, किसी भी अतिरिक्त वर्तमान प्रोटोकॉल शुल्क और किसी भी अलग से प्रस्तावित रूटिंग या सेवा शुल्क को सूचीबद्ध किया जाएगा। निष्पादन समाप्त हो चुके उद्धरणों या उल्लंघन सीमाओं को अस्वीकार कर देगा।

\textbf{परमाणु शक्ति और पुनर्प्राप्ति।} मल्टी-हॉप निष्पादन जहां संभव हो परमाणु होगा। जब यह एचटीएलसी या एस्क्रो का उपयोग करता है, तो सेवा टाइमआउट, अवरोध पथ, जिम्मेदार नियंत्रकों, घटना प्रक्रियाओं और शेष जोखिमों का खुलासा करेगी।

\textbf{प्रस्तावित बैच नेटिंग और पुनः संतुलन।} एक ऑप्ट-इन सेवा संतुलन पुनर्प्राप्त करने के इरादों को इकट्ठा कर सकती है और संगत चक्र या श्रृंखलाओं की खोज कर सकती हैः

\begin{enumerate}[leftmargin=*]
\item निष्पादित चक्रों, परिसंपत्तियों, मात्राओं, मूल्य निर्धारण समयशीर्षक और शुल्क की पहचान करने वाली एक मशीन-पठनीय रसीद प्रकाशित करें;
\item अपनाई गई प्रति अवधि सीमाओं और प्रतिपक्ष नीति को लागू करना;
\item किसी भी भाग लेने वाले पूल के प्राधिकरण, सीमाओं या उपलब्ध सूची का उल्लंघन करने वाली गतिविधि को अस्वीकार करना; तथा
\item समीक्षा और विवाद निपटान के लिए निर्धारक इनपुट और प्राप्तियों को संरक्षित करना।
\end{enumerate}
\textbf{SDK आवश्यकताएँ।} एकSDKनिष्पादित मार्गों के लिए निर्धारक उद्धरण-से-प्राप्ति मानचित्रण, प्रति हॉप अपरिवर्तित जांच, समझने योग्य विफलता कोड और लेखा परीक्षा अनुकूल लॉग प्रदान करेगा।Protocol v1.1.0 \texttt{SwapRouter} केवल उद्धरण प्रदान करता है; यह इन प्रस्तावित मार्गों को निष्पादित नहीं करता है।

\subsubsection*{8.1.1 न्यूनतम संघ संगतता विनिर्देश}

एक पूल पारिस्थितिकी तंत्र जो क्रॉस-प्रोफाइल राउटिंग की मांग करता है, मशीन-पठनीय जानकारी प्रकाशित करेगाः

\begin{enumerate}[leftmargin=*]
\item \textbf{रजिस्ट्री जड़ेंः} परिसंपत्तियों, पूल, विनिमय दर विधियों, सीमाओं और शुल्क नीतियों के लिए पहचानकर्ता, या एक रूट जो उन्हें निर्णायक रूप से हल करता है।
\item \textbf{रसीदेंः} प्रोफ़ाइल, अंदर और बाहर की संपत्ति, राशि, उद्धरण स्रोत और समय टिकट, सीमा स्नैपशॉट, शुल्क, इन्वेंट्री परिणाम और प्रत्येक छलांग के लिए निष्पादन परिणाम।
\item \textbf{ऑपरेशनल सिग्नलः} इन्वेंट्री, सीमित उपयोग, घटनाओं और किसी भी अलग से प्रमाणित वादे की पूर्ति या वित्त पोषित संरक्षण के बारे में ताजाता-सीमित जानकारी।
\item \textbf{नीतिगत प्रतिबंधः} अनुमत या अस्वीकार किए गए प्रतिद्वंद्वी, परिसंपत्ति वर्ग, एडाप्टर और किसी भी एस्क्रो आवश्यकताएं।
\item \textbf{विफलता कोडः} अस्वीकृति, समाप्ति, सीमा, इन्वेंट्री, नीति, निर्भरता या घटना विफलताओं के लिए निर्णायक स्पष्टीकरण।
\end{enumerate}
एक प्रोफ़ाइल बुनियादी CPP संगतता के लिए आवश्यकताओं के बिना कवरेज, अनुपालन, मध्यस्थता या अन्य सेवाओं को जोड़ सकता है। प्रत्येक वैकल्पिक सेवा अपने जिम्मेदार पक्ष, प्राधिकरण, दायरे और शर्तों की पहचान करेगी।

\subsection*{8.2 लाइसेंस, सत्यापन और बाहर निकलना}
\addcontentsline{toc}{subsection}{8.2 लाइसेंस, सत्यापन और बाहर निकलना}

Protocol v1.1.0 अनुबंधEVM-प्रोटोकॉल रिपॉजिटरी में अनुबंध \texttt{src} निर्देशिका के अंतर्गत प्रकाशित की जाती हैAGPL-3.0अपवाद के लिए पहचाने गए अपरिवर्तित तृतीय-पक्ष घटक जो अपनी शर्तों को बनाए रखते हैं। प्रकाशित स्रोत, एबीआई और तैनाती निर्देश स्वतंत्र समीक्षा का समर्थन करते हैं लेकिन स्वयं एक लेखा परीक्षा, सुरक्षित तैनाति या कानूनी अनुपालन साबित नहीं करते हैं।

प्रत्येक तैनाती अपने कोड संस्करण, बिल्ड मूल, पते, नियंत्रक और अपग्रेड शक्तियां, ऑडिट स्थिति, रजिस्ट्री दर्पण और किसी भी टाइमलॉक या प्यूज सुरक्षा को अलग से प्रकट करेगी।

एक प्रस्तावित \textbf{फोर्क किट} इसमें शामिल हो सकते हैंः

\begin{enumerate}[leftmargin=*]
\item निर्धारक तैनाती स्क्रिप्ट;
\item रजिस्टर स्नैपशॉट और निर्यात उपकरण;
\item मार्ग सेवाओं, एसडीके और इंटरफेस को एक नई रजिस्ट्री रूट पर पुनर्निर्देशित करने के लिए एक प्रलेखित प्रक्रिया;
\item साझा रजिस्टर को सुरक्षित रूप से छोड़ने के लिए पूल प्रबंधक चेकलिस्ट; तथा
\item बकाया vouchers के लिए एक प्रवास चेकलिस्ट, जिसमें जारीकर्ता सूचनाएं, प्रस्तुत करने और पूरा करने की समय सीमा, रिकॉर्ड तक निरंतर पहुंच और सुधार शामिल हैं।
\end{enumerate}
संगत कांटे तब लचीलापन में सुधार कर सकते हैं जब समुदायों, सहकारी संस्थाओं, सार्वजनिक एजेंसियों, संघों, मल्टी-सिग्स या सेवा ऑपरेटरों को अलग शासन की आवश्यकता होती है। वास्तविक निरंतरता अभी भी अनुबंध स्वामित्व, कुंजी, निर्भरताओं, इंटरफ़ेस, बुनियादी ढांचे, कानूनी दायित्वों और तृतीय पक्ष सेवाओं पर निर्भर करती है।


\section*{9. तरलता कार्यक्रमों के लिए प्रस्तावित अर्थशास्त्र}
\addcontentsline{toc}{section}{9. तरलता कार्यक्रमों के लिए प्रस्तावित अर्थशास्त्र}

इस अध्याय में एक भविष्य के मॉडल का वर्णन किया गया है, जिसे अलग से अपनाया गया है। यह एक वर्तमान ऐप सुविधा, एक प्रस्ताव, एक वादा रिटर्न या \texttt{SwapPool} में जमा करके बनाया गया अधिकार नहीं है।

\subsection*{9.1 प्रस्तावित राजस्व स्रोत}
\addcontentsline{toc}{subsection}{9.1 प्रस्तावित राजस्व स्रोत}

भविष्य के नेटवर्क बजट में निम्नलिखित राशि दी जा सकती हैः

\begin{enumerate}[leftmargin=*]
\item एक खुलासा \textbf{नेटवर्क रैक} भाग लेने वाले पूल द्वारा एकत्रित शुल्क के हिस्से के रूप में लिया गया;
\item लागू साझा सेवाओं से अलग रूटिंग या सेवा शुल्क; तथा
\item अन्य स्पष्ट रूप से प्राप्त, प्राप्त राजस्व।
\end{enumerate}
पूल द्वारा बनाए गए सकल पूल शुल्क नेटवर्क राजस्व नहीं हैं। वर्तमान Protocol v1.1.0 एक अलग मॉडल का उपयोग करता हैः एक वैकल्पिक प्रोटोकॉल शुल्क पूल फीस के अतिरिक्त है और सीधे इसके कॉन्फ़िगर किए गए प्राप्तकर्ता को भेजा जाता है।

\subsection*{9.2 चित्रणात्मक रैक गणित}
\addcontentsline{toc}{subsection}{9.2 चित्रणात्मक रैक गणित}

यदि एक भाग लेने वाले पूल से 2.00\% पूल शुल्क लिया जाता है और एक अपनाए गए नेटवर्क रैक को उस पूल फीस का 20\% प्राप्त होता है, तो रूट किए गए मूल्य पर प्रस्तावित प्रभावी नेटवर्क-रेक दर हैः

\texttt{rake\_rate = 2.00\% $\times$ 20\% = 0.40\% = 40 bps}

यदि प्राप्त रैक और सेवा शुल्क परिसंपत्तियों का 25\% लागत के बाद घोषित नकदी में नामित उपयोग के लिए पात्र और परिवर्तनीय है, तो 40 बीपीएस रैक का नकद-उपयोग योग्य हिस्सा लगभग 10 बीपीसी है।

प्रस्तावित नेटवर्क राजस्व:

\texttt{network\_rake\_received + routing\_or\_service\_fees\_received}

नेटवर्क रैक में सकल पूल शुल्क न जोड़ेंः रैक उन फीस से एक हस्तांतरण है और अन्यथा दो बार गिना जाएगा।

\subsection*{9.3 तरलता कार्यक्रम अधिकार}
\addcontentsline{toc}{subsection}{9.3 तरलता कार्यक्रम अधिकार}

एक अलग से अपनाया गया कार्यक्रम पूल इन्वेंट्री, रूटिंग सेवाओं, निगरानी या अन्य जनादेशों को वित्त पोषित कर सकता है। इसकी शर्तों का खुलासा करने की आवश्यकता होगीः

\begin{itemize}[leftmargin=*]
\item क्या एक हस्तांतरण उपहार, दान, ऋण, वसूली योग्य योगदान या खरीद है;
\item रखरखाव और नियंत्रण;
\item निकासी, वापसी, हानि और प्राथमिकता के नियम;
\item शुल्क और इनाम की पात्रता;
\item शासन अधिकार;
\item मूल्यांकन और रिपोर्टिंग विधियां; तथा
\item निलंबन, समाप्ति और उपचार।
\end{itemize}
वर्तमान \texttt{SwapPool} में कोई पूल-शेयर टोकन या स्वचालित योगदानकर्ता अधिकार नहीं होता है। किसी भी एक्स-पोस्ट मीट्रिक को वास्तविक प्राप्तियों और नुकसान पर आधारित होना चाहिए, वादा किए गए रिटर्न के रूप में प्रस्तुत नहीं किया जाना चाहिए, और शून्य या नकारात्मक हो सकता है।


\section*{10. व्यापक जोखिम ढांचा}
\addcontentsline{toc}{section}{10. व्यापक जोखिम ढांचा}

इस प्रस्तावित ढांचे में दस जोखिम श्रेणियों को अलग किया गया है। प्रत्येक श्रेणी के लिए यह संभावित संकेतकों, नियंत्रणों, तनाव परीक्षणों और एक दिशानिर्देशात्मक जोखिम भूख की पहचान करता है। ये डिजाइन सिफारिशें हैं, यह दावा नहीं करते कि हर नियंत्रण तैनात, प्रभावी या पर्याप्त है। सीमाएं, भंडार, गारंटी, निगरानी, कवरेज और शासन हानि को समाप्त नहीं कर सकते।

\subsection*{10.1 प्रोटोकॉल और स्मार्ट कॉन्ट्रैक्ट जोखिम}
\addcontentsline{toc}{subsection}{10.1 प्रोटोकॉल और स्मार्ट कॉन्ट्रैक्ट जोखिम}

\begin{itemize}[leftmargin=*]
\item \textbf{खतरे:} अनुबंध कीड़े, उन्नयन त्रुटियों, निर्भरता विफलताओं और गलत रूप से कॉन्फ़िगर सीमाओं या शुल्क।
\item \textbf{संकेतक:} लेखा परीक्षा निष्कर्ष, अनौपचारिक इन्वेंट्री आंदोलन, अपरिवर्तित विफलताएं और असामान्य रिवर्स।
\item \textbf{संभव नियंत्रण:} स्वतंत्र लेखा परीक्षाएं; न्यूनतम विशेषाधिकारित भूमिकाएँ; खुलासा किए गए प्रॉक्सी और निर्भरता नियंत्रक; समय में लॉक अपग्रेड; श्रृंखला पर निगरानी; प्रकाशित प्राधिकरण और मानदंडों के साथ घटना रुकावटें; और एक परीक्षण किया गया प्रवासन पथ।
\item \textbf{तनाव परीक्षण:} अनुपलब्ध बोली या सीमा निर्भरता, इन्वेंट्री की कमी, ठहराव वाले अनुबंध और ट्रैफ़िक फट गया।
\item \textbf{जोखिम की इच्छा:} बकाया दायित्वों या स्वैप मात्रा को स्केल करने से पहले कम।
\end{itemize}
\subsection*{10.2 आर्थिक और बाजार जोखिम}
\addcontentsline{toc}{subsection}{10.2 आर्थिक और बाजार जोखिम}

\begin{itemize}[leftmargin=*]
\item \textbf{खतरे:} बारीक इन्वेंट्री, एकतरफा प्रवाह, तेजी से निकासी और हेरफेर किए गए या पुराने मूल्य संदर्भ।
\item \textbf{संकेतक:} उच्च पूल टोकन बैलेंस कैप उपयोग, बढ़ते उद्धरण अंतर, बार-बार सीमा अस्वीकार और केंद्रित इन्वेंट्री।
\item \textbf{संभव नियंत्रण:} वर्तमान पूल टोकन बैलेंस कैप; प्रस्तावित रोलिंग या अकाउंट लिमिट; अलग से अपनाए गए रिजर्व; संरक्षित मूल्य संदर्भ; मार्ग बहिष्करण; और समय-सीमित घटना शुल्क या सीमाएं।
\item \textbf{तनाव परीक्षण:} बड़े संदर्भ मूल्य आंदोलन, प्रस्तुति वृद्धि, निकासी के अनुरोध और डेटा स्रोतों में विराम।
\item \textbf{जोखिम की इच्छा:} केवल प्रकाशित मापदंडों और वित्त पोषित नुकसान सहन क्षमता के भीतर मध्यम।
\end{itemize}
\subsection*{10.3 जारीकर्ता और वाउचर जोखिम}
\addcontentsline{toc}{subsection}{10.3 जारीकर्ता और वाउचर जोखिम}

\begin{itemize}[leftmargin=*]
\item \textbf{खतरे:} निष्पादन क्षमता से परे जारी, emitter default, misleading terms और poorly specified availability or presentation windows।
\item \textbf{संकेतक:} पुष्टि की गई वादे की पूर्ति दरों में गिरावट, बुढ़ापे वाले बकाया वाउचर, अनसुलझे शिकायतें और एक जारीकर्ता के प्रति केंद्रित जोखिम।
\item \textbf{संभव नियंत्रण:} issuer due diligence; clear voucher and offer terms; issuance or admission limits; independently funded bonds or guarantees; cohort reporting; और accountable registry review.
\item \textbf{तनाव परीक्षण:} जारीकर्ता की दिवालियापन, क्षेत्रीय उत्पादन झटके, नकली दावे और लंबे समय तक पूरा होने में देरी।
\item \textbf{जोखिम की इच्छा:} जारीकर्ता या प्रस्ताव की एकाग्रता में वृद्धि के साथ कम।
\end{itemize}
\subsection*{10.4 मुआवजा प्रस्तुत करने और वादे की पूर्ति का जोखिम}
\addcontentsline{toc}{subsection}{10.4 मुआवजा प्रस्तुत करने और वादे की पूर्ति का जोखिम}

\begin{itemize}[leftmargin=*]
\item \textbf{खतरे:} अवैध या डुप्लिकेट प्रस्तुति, निष्पादक की अपर्याप्त क्षमता, स्टॉकआउट, रसद विफलता और अनुपलब्ध डिस्चार्ज रिकॉर्ड।
\item \textbf{संकेतक:} अनुरोध पूरा करने में देरी, असफल या विवादास्पद प्रस्तुतियाँ, स्टॉकआउट, टिकट बैकलॉग और पूर्ण इकाइयां जिन्हें डिस्चार्ज नहीं मिला है।
\item \textbf{संभव नियंत्रण:} प्रकाशित प्रस्तुति और वादे की पूर्ति प्रक्रियाएं; क्षमता का खुलासा; वैध होने पर कई वादे की पूर्ति स्थल; साक्ष्य मानक; शिकायत और उपचार पथ; और प्रस्तुति, वादे की पूर्ति और डिस्चार्ज को अलग करने वाले रिकॉर्ड।
\item \textbf{तनाव परीक्षण:} दो से चार गुना प्रस्तुति मात्रा, सुविधाओं के टूटने, आपूर्तिकर्ता की विफलता और दोहरा उपयोग प्रयास।
\item \textbf{जोखिम की इच्छा:} कम जहां आवश्यक वस्तुओं, कमजोर प्रतिभागियों या लंबी वादे की पूर्ति खिड़कियां शामिल हैं।
\end{itemize}
\subsection*{10.5 शासन जोखिम}
\addcontentsline{toc}{subsection}{10.5 शासन जोखिम}

\begin{itemize}[leftmargin=*]
\item \textbf{खतरे:} Steward या नियंत्रक कैप्चर, त्वरित पैरामीटर परिवर्तन, हितों के संघर्ष, छिपी तकनीकी शक्तियों, और कमजोर अपील।
\item \textbf{संकेतक:} एकाग्रता, आवधिक आपातकालीन कार्रवाई, नीतिगत परिवर्तन और दोहराए जाने वाले ओवरराइड।
\item \textbf{संभव नियंत्रण:} भूमिका और शक्ति का खुलासा; समानुपातिक अनुमोदन की सीमाएँ; समय-सीमाएं; संघर्ष और अस्वीकृति नियम; सार्वजनिक परिवर्तन रिकॉर्ड; अपीलें; स्वचालित आपातकालीन बिजली सूर्यास्त; और फोर्कबिलिटी।
\item \textbf{तनाव परीक्षण:} विपक्षी प्रस्ताव, हस्ताक्षरकर्ता हानि, रिश्वत के प्रयास और संचय या अधिकृत नियंत्रण के माध्यम से कब्जा।
\item \textbf{जोखिम की इच्छा:} मूल्य पद्धति, निकासी, रजिस्ट्री रूट, कवरेज या आपातकालीन शक्तियों को प्रभावित करने वाले कार्यों के लिए कम।
\end{itemize}
प्रस्तावित CLC शासन टोकन के भविष्य की तैनाती के लिए, कैप्चर परीक्षण में बजट को पुनर्निर्देशित करने, रजिस्ट्री मानकों को कमजोर करने, संबंधित पक्षों के जनादेशों को मंजूरी देने या वित्त पोषित कवरेज को समाप्त करने के प्रयासों के बाद संचय शामिल होगा। संभावित सुरक्षा उपायों में शासन लॉकअप, महत्वपूर्ण कार्यों के लिए उच्च सीमाएं, देरी से निष्पादन, पारदर्शी प्रतिनिधि और एकाग्रता निगरानी, एक घटना प्रक्रिया और एक विश्वसनीय फोर्क-एंड-एक्सिट प्रक्रिया शामिल है। कोई भी यहां दिखाई देकर ही तैनात के रूप में प्रतिनिधित्व नहीं किया गया है।

\subsection*{10.6 कानूनी और अनुपालन जोखिम}
\addcontentsline{toc}{subsection}{10.6 कानूनी और अनुपालन जोखिम}

\begin{itemize}[leftmargin=*]
\item \textbf{खतरे:} एक वाउचर, सेवा, पदोन्नति या शासन संपत्ति जो अप्रत्याशित नियामक उपचार प्राप्त करती है; उपभोक्ता संरक्षण में विफलता; धन शोधन या प्रतिबंधों का जोखिम; और सीमा पार प्रतिबंध।
\item \textbf{संकेतक:} अधिकार क्षेत्र के ध्वज, नियामक पूछताछ, शिकायतें, प्रतिबंधित पक्षों के मैच और विज्ञापन और वास्तविक व्यवहार के बीच विचलन।
\item \textbf{संभव नियंत्रण:} श्रेणी और अधिकार क्षेत्र द्वारा समीक्षा; सटीक प्रकटीकरण; भू-आधारित इंटरफेस; समानुपातिक पात्रता या प्रमाणन जांच; पदोन्नति नियंत्रण; प्राधिकरण और स्वीकृति के रिकॉर्ड; और स्पष्ट जिम्मेदार पक्ष।
\item \textbf{तनाव परीक्षण:} एक क्षेत्राधिकार प्रतिबंध, प्रदाता समाप्ति, अनिवार्य पुनः वर्गीकरण और किसी विशेषता या परिसंपत्ति वर्ग को रोकने का आदेश।
\item \textbf{जोखिम की इच्छा:} कम; असमर्थित गतिविधि को प्रतिबंधित या रोकें।
\end{itemize}
\subsubsection*{10.6.1 कानूनी स्थिति और प्रस्तावित टोकन उपचार}

\begin{enumerate}[leftmargin=*]
\item \textbf{सत्यापित बुनियादी ढांचाः} Protocol v1.1.0 अनुबंधEVM-प्रोटोकॉल रिपॉजिटरी में पहचान किए गए लाइसेंस और तीसरे पक्ष के अपवादों के तहत प्रकाशित किया गया है। प्रकाशन समीक्षा की अनुमति देता है लेकिन स्वयं एक लेखा परीक्षा, सुरक्षित तैनाती या कानूनी अनुपालन का प्रमाण नहीं देता है। प्रत्येक तैनाति को अपने कोड संस्करण, निर्माण मूल, पते, नियंत्रक शक्तियों और लेखा परीक्षा स्थिति का खुलासा करना चाहिए।
\item \textbf{प्रस्तावित टोकन मुद्राः} इस डिजाइन में, प्रस्तावित CLC शासन टोकन शासन और नीतिगत प्रवेश को समन्वयित करेगा। यह लाभांश, लाभ साझाकरण, अवशिष्ट अधिकार या मूल्य या तरलता की गारंटी नहीं देगा।
\item \textbf{संबंधित प्रस्तावित परिसंपत्तियांः} एक अलग से लागू नीति प्रस्तावित लॉक करने की अनुमति दे सकती हैCLCशासन का प्रतीकstCLCऔर युग-अवसर के लिए अधिकृत कर सकते हैंsCLC. अपनाए गए शर्तों में उनके अधिकार, सीमाएं, समाप्ति, हस्तांतरण और उपचार को सटीक रूप से परिभाषित करना होगा।
\item \textbf{संचार:} किसी भी प्रस्तावित सामग्री के लिएCLCशासन-सकेत,stCLCयाsCLCतैनाती को लाभ, मूल्य निर्धारण, निष्क्रिय आय या गारंटीकृत पहुंच का वादा नहीं करना चाहिए।
\item \textbf{अधिकार क्षेत्र नियंत्रण:} एक तैनाती के लिए इंटरफेस जियोफेंसिंग, प्रतिबंधित वर्गों के लिए प्रमाणपत्र, पदोन्नति सीमाओं, परिसंपत्ति-विशिष्ट समीक्षा या नियंत्रण की आवश्यकता हो सकती है जो प्रस्तावित सुविधा को रोकते हैं।
\item \textbf{प्रतिभागियों की सूचनाः} स्वीकृत शर्तें बताएंगी कि कानूनी, परिचालन या जोखिम कारणों से पहुंच को कब कम या अक्षम किया जा सकता है और क्या कोई मुआवजा या सुधार लागू होता है।
\end{enumerate}
\subsection*{10.7 रूटिंग और क्रॉस-डोमेन जोखिम}
\addcontentsline{toc}{subsection}{10.7 रूटिंग और क्रॉस-डोमेन जोखिम}

\begin{itemize}[leftmargin=*]
\item \textbf{खतरे:} आंशिक निष्पादन, फंसे हुए हॉप्स, ब्रिज या एस्क्रो शोषण, पुराने कोटेशन, पथ मुद्रास्फीति और घोषित मूल्य परिवर्तनों का अग्रिम संचालन।
\item \textbf{संकेतक:} मार्ग की समाप्ति दरें, एस्क्रो बैकलॉग्स, कोट-टू-इक्जीक्यूशन अंतर, दोहराए जाने वाले अनावश्यक कूद और पुल घटनाएं।
\item \textbf{संभव नियंत्रण:} यदि उपलब्ध हो तो परमाणु निष्पादन; रूढ़िवादी HTLC या एस्क्रो टाइमआउट; पथ और प्रतिपक्ष नीति; मार्ग स्तर की सीमाएं; निर्णायक बोली-से-प्राप्ति मानचित्रण; और उत्तरदायी सेवा ऑपरेटर।
\item \textbf{तनाव परीक्षण:} एक पुल बंद, श्रृंखला पुनर्गठन, निर्भरता टूटने, और प्रस्तावित मल्टी-हॉप मार्ग में एक असफल कूद।
\item \textbf{जोखिम की इच्छा:} केवल पहचान और निगरानी की गई निर्भरता के लिए कम से मध्यम।
\end{itemize}
\subsection*{10.8 रखरखाव और कुंजी प्रबंधन का जोखिम}
\addcontentsline{toc}{subsection}{10.8 रखरखाव और कुंजी प्रबंधन का जोखिम}

\begin{itemize}[leftmargin=*]
\item \textbf{खतरे:} खोए हुए या बाधित कुंजी, हस्ताक्षरकर्ता संदिग्धता और अस्पष्ट पुनर्प्राप्ति या प्रशासक प्राधिकरण।
\item \textbf{संकेतक:} अनियमित हस्ताक्षर, नियंत्रक परिवर्तन, विफल रोटेशन और असामान्य निकासी।
\item \textbf{संभव नियंत्रण:} मल्टीसिग या थ्रेसहोल्ड प्राधिकरण; हार्डवेयर समर्थित कुंजी; भूमिका पृथक्करण; हस्ताक्षरकर्ता रोटेशन; सार्वजनिक नियंत्रक सूची; निगरानी किए गए निकासी सीमाएं; और परीक्षण की गई वसूली प्रक्रियाएँ।
\item \textbf{जोखिम की इच्छा:} कम।
\end{itemize}
\subsection*{10.9 प्रतिष्ठा और सामाजिक जोखिम}
\addcontentsline{toc}{subsection}{10.9 प्रतिष्ठा और सामाजिक जोखिम}

\begin{itemize}[leftmargin=*]
\item \textbf{खतरे:} भ्रामक दावे, हानिकारक प्रोत्साहन, inaccessible complaints, poor fulfillment experiences, privacy failures, and protections that favor insiders।
\item \textbf{संकेतक:} समूहों द्वारा शिकायतें, अनसुलझे विवाद, जारीकर्ता-प्रदर्शन के रुझान, लाभ या हानि का एकाग्रता और सामुदायिक प्रतिक्रिया।
\item \textbf{संभव नियंत्रण:} स्पष्ट भाषा में प्रकटीकरण; शिकायत और सुधार पथ; सुलभ साक्ष्य; पारदर्शी रिपोर्टिंग; घटनाओं की समीक्षा; और गलत बयान के लिए समानुपातिक दंड।
\item \textbf{जोखिम की इच्छा:} कम, विशेष रूप से प्रभावित समुदायों और कमजोर प्रतिभागियों की देखभाल के साथ।
\end{itemize}
\subsection*{10.10 एकाग्रता और विखंडन का जोखिम}
\addcontentsline{toc}{subsection}{10.10 एकाग्रता और विखंडन का जोखिम}

\begin{itemize}[leftmargin=*]
\item \textbf{खतरे:} कम संख्या में जारीकर्ताओं, पूल, नियंत्रकों, प्रदाताओं, नेटवर्क या असंगत कांटे पर निर्भरता।
\item \textbf{संकेतक:} जारीकर्ता, पूल, इन्वेंट्री, नियंत्रक या सेवा प्रदाता द्वारा एकाग्रता उपाय; समूहों के बीच मार्ग विफलताएं; और महत्वपूर्ण एकल निर्भरताएँ।
\item \textbf{संभव नियंत्रण:} प्रकाशित एकाग्रता सीमाएं; कई उत्तरदायी ऑपरेटर; संगत मानक; स्वतंत्र रजिस्टर; परीक्षण किए गए प्रस्थान प्रक्रियाएँ; और सुरक्षित सहकार्य।
\item \textbf{तनाव परीक्षण:} सबसे बड़ा जारीकर्ता, पूल, ऑपरेटर, प्रदाता या रजिस्ट्री रूट का नुकसान।
\item \textbf{जोखिम की इच्छा:} आवश्यक सेवाओं या अपरिवर्तनीय निर्भरताओं के लिए अधिक सख्त सीमाओं के साथ तैनाती-विशिष्ट और खुलासा।
\end{itemize}

\section*{11. शासन तंत्र}
\addcontentsline{toc}{section}{11. शासन तंत्र}

इस अध्याय में एक शासन टेम्पलेट प्रस्तावित किया गया है। यह प्रतिनिधित्व नहीं करता है कि वर्तमान CLC App शासन टोकन मतदान, टाइमलॉक, साझा बीमा, दावा प्रक्रिया या नीचे वर्णित सभी नियंत्रण का उपयोग करता है। प्रत्येक तैनाती को अपने वास्तविक निर्णय निर्माताओं, अधिकारियों, अनुबंधों, प्रक्रियाओं और नीतियों की पहचान करने की आवश्यकता होगी।

\begin{itemize}[leftmargin=*]
\item \textbf{संवैधानिक मूल्य:} लोगों के लिए देखभाल, पर्यावरण की देखभाल, निष्पक्षता, पारस्परिकता, गैर-प्रभुत्व और लचीलापन।
\item \textbf{प्रस्तावों के प्रकारः} शुल्क, सीमा और सूचकांक परिवर्तन; तरलता जनादेश; पूल लिस्टिंग और निकासी; वैकल्पिक कवरेज निर्णय; और पैरामीटर गार्डरेल्स।
\item \textbf{उत्तरदायी प्रक्रियाः} उपभोग $\to$ मूल्यांकन $\to$ जोखिम समीक्षा $\to$ अनुमोदन $\to$ समय सीमा जहां उपयुक्त हो $\to$ निष्पादन। अनुमोचन प्रबंधक, सहकारी संस्थाओं, सार्वजनिक एजेंसियों, संघों, मल्टी-सिग्स, श्रृंखला पर मतदान या किसी अन्य खुलासा और जवाबदेही संरचना से आ सकता है।
\item \textbf{अनुमोदन की सीमाएंः} मूल्य सूचकांक परिवर्तनों, आपातकालीन शक्तियों और अन्य महत्वपूर्ण कार्यों के लिए उच्चतम सीमाओं के साथ कार्रवाई वर्ग द्वारा मापदंडित।
\item \textbf{प्रतिनिधिमंडल:} सार्वजनिक जनादेशों के साथ वैकल्पिक प्रतिनियुक्ति, संघर्ष का खुलासा और वापसी।
\item \textbf{सर्किट ब्रेकरः} निर्दिष्ट मानदंडों के साथ आपातकालीन विराम, अधिकृत ऑपरेटर, फिर से शुरू करने की शर्तें और आवश्यक पोस्ट-मॉर्टम।
\item \textbf{पारदर्शिताः} स्वैप सेटलमेंट, जारीकर्ता वादे की पूर्ति, भंडार, लिमिट उपयोगिता, रूटिंग और गारंटर के लिए अलग सबूतों के साथ प्रकाशित परिवर्तन और प्रवाह।
\end{itemize}
\textbf{रजिस्टर शासन।} एक CPP- संगत तैनाती वाउचर, टोकन और पूल के लिए डिस्कवरी रजिस्टर बनाए रख सकती है। अधिकृत नियंत्रण तैनाति की खुलासा शासन प्रक्रिया के माध्यम से रजिस्ट्री प्रविष्टियों को जोड़, अद्यतन, निलंबित या हटा सकते हैं। रजिस्ट्रार हटाने से उस रजिस्ट्रे में खोज और रूटिंग प्रभावित होती है; यह अपने आप में एक टोकन को नहीं मिटाता है, धारक का बैलेंस नहीं बदलता, जारीकर्ता के दायित्व को समाप्त करता है, या अन्यथा कार्यात्मक अनुबंध को अक्षम करता है।

प्रकाशित रजिस्टर नियमों में स्थिति को सशर्त बनाना चाहिए और इसे निलंबन या हटाने के लिए आधार के रूप में दोहराए गए गैर-वादे की पूर्ति, धोखाधड़ी या गलत प्रतिनिधित्व, असुरक्षित अनुबंध व्यवहार, या प्रकाशित सिद्धांतों का लगातार उल्लंघन की पहचान कर सकता है। जहां संभव हो, प्रक्रिया को सूचना, सुधार का अवसर और अपील पथ प्रदान करना चाहिए। आपातकालीन हटाने से सार्वजनिक घटना रिपोर्ट और स्वचालित समीक्षा या सूर्यास्त की आवश्यकता होनी चाहिए।

\textbf{प्रतिबंधित सूची।} इस टेम्पलेट के तहत, एक रजिस्टर स्वीकार नहीं करेगाः

\begin{enumerate}[leftmargin=*]
\item ऐसे उपकरण जो सीधे सहमत सीमाओं से परे पारिस्थितिक विनाश, हिंसा या सशस्त्रीकरण, जबरन निष्कर्षण या प्रणालीगत दुरुपयोग को वित्त पोषित या प्रोत्साहित करते हैं; या
\item एक वाउचर वर्ग जिसमें स्पष्ट प्रस्तुति और वादे की पूर्ति की शर्तें, जवाबदेही और उपचार मार्गों का अभाव है।
\end{enumerate}
प्रतिबंधित सूची को केवल अनुमोदित महत्वपूर्ण कार्रवाई की सीमा और समय सीमा के माध्यम से संस्करण, सार्वजनिक रूप से लेखा परीक्षा और परिवर्तन योग्य किया जाएगा, जैसा कि चित्रित हैQ3 + T3परिशिष्ट D में।

\subsection*{11.1 विनिमय दर और सीमा शासन}
\addcontentsline{toc}{subsection}{11.1 विनिमय दर और सीमा शासन}

\textbf{समय पर लॉक परिवर्तन।} इस टेम्पलेट का अनुसरण करने वाली एक तैनाती केवल सार्वजनिक समय लॉक के बाद विनिमय दर विधियों को बदल देगी और सीमा मापदंडों को अलग से लागू करेगी। आपातकालीन पथ में एक अलग से खुलासा की गई प्राधिकरण प्रक्रिया का उपयोग किया जाएगा और इसमें स्वचालित सूर्यास्त या समीक्षा शामिल होगी।

\textbf{अनुमोदन की सीमाएं।} इस टेम्पलेट में मूल्य सूचकांक आधार परिवर्तनों और वैश्विक सीमा स्तर परिवर्तनों, पूल-विशिष्ट तृतीय-पक्ष परिवर्तनों के लिए मध्यवर्ती सीमाओं और नियमित शुल्क परिवर्तन के लिए मानक सीमाओं के लिए उच्च अनुमोदन दरें प्रस्तावित हैं।

\textbf{प्रकाशित फ़ीड.} एक भाग लेने वाली तैनाती प्रत्येक पूल के लिए ऑन-चेन सूचकांक चर, ओरेकल स्रोतों या माध्यमों, अद्यतन कैडेन्स, सीमा खिड़कियों और टोपी, और विफलता मोड या सुरक्षित स्थिरांक प्रकाशित करेगी।

\textbf{आपातकालीन विराम के मानदंड।} एक भाग लेने वाली तैनाती में पूर्व-घोषणा की जाएगी जैसे कि ओरेकल आउटेज, पूर्णता विफलताओं के साथ उच्च सीमा उपयोग या अपरिवर्तनीय विफलता, साथ ही फिर से शुरू जांच और घटना के बाद समीक्षा आवश्यकताओं।

\subsection*{एक पूल और वाउचर के लिए सार्वजनिक सूचकांक फ़ीड का उदाहरण}
\addcontentsline{toc}{subsection}{एक पूल और वाउचर के लिए सार्वजनिक सूचकांक फ़ीड का उदाहरण}

\begin{itemize}[leftmargin=*]
\item \textbf{प्रतीक:} उदाहरण के लिए \texttt{Maize\_50kg@IssuerY}।
\item \textbf{संदर्भ इकाईः} सूचकांक इकाई (IUX)
\item \textbf{प्रकाशित मूल्यः} 30.000 IUX.
\item \textbf{स्रोत:} स्थानीय बाजार सर्वेक्षण, मंत्रालय की बुलेटिन और तैनाती आधार रेखा जैसे ज्ञात स्रोतों का औसत।
\item \textbf{अद्यतन क्रमशः} प्रतिदिन 18:00 बजे EAT, 24 घंटे के समय लॉक के साथ।
\item \textbf{विफलता मोडः} अंतिम वैध मूल्य पर फ्रीज करें, एक खुलासा सीमा नीति लागू करें और 72 घंटे के विराम के बाद रुकें।
\item \textbf{तर्कः} प्रकाशित नोट्स और पिछले अद्यतन से परिवर्तन रिकॉर्ड।
\item \textbf{हस्ताक्षरकर्ता:} खुलासा किए गए मल्टी सिग पते और अनुमोदन की सीमा।
\end{itemize}
\subsection*{11.2 प्रस्तावित बीमा निधि रनबुक}
\addcontentsline{toc}{subsection}{11.2 प्रस्तावित बीमा निधि रनबुक}

\textbf{केवल वैकल्पिक डिजाइन।} यह रनबुक केवल एक तैनाती पर लागू होता है जिसने स्पष्ट रूप से बीमा कोष को अपनाया और वित्त पोषित किया है और कवर की गई घटनाओं, पात्र दावेदारों, जिम्मेदार इकाई, संपत्ति, सीमाएं, बहिष्करण, साक्ष्य आवश्यकताओं, प्रक्रिया और शासन शर्तों को प्रकाशित किया है।CLC App और न हीGEF यह केवल इसलिए है क्योंकि यह डिजाइन श्वेत पत्र में दिखाई देता है।

\textbf{संभावित ट्रिगर।} एक स्वीकृत नीति परिभाषित जारीकर्ता अनुपालन, पूल रिजर्व घाटा, या पुल या एस्क्रो हानि को कवर कर सकती है। एक तकनीकी घटना स्वचालित रूप से योग्य नहीं होती है; लागू नीति नियंत्रण करेगी।

\textbf{मूल्यांकन।} जवाबदेह निकाय transaction receipts, inventory balances, guarantor bonds, भुनाने की प्रस्तुति के records, जारीकर्ता के जवाब और दूसरे ज़रूरी प्रमाणों को मिलाएगा, फिर privacy और कानून के अनुसार incident record प्रकाशित करेगा।

\textbf{स्पष्ट रूप से नुकसान झरना।} जहां प्रत्येक परत मौजूद है और कानूनी रूप से लागू होती है, तो एक पॉलिसी का उपयोग किया जा सकता हैः (1) उत्तरदायी जारीकर्ता बांड या गारंटर दांव $\to$ (2) पूल स्तर के भंडार $\to$ (3) एक प्रस्तावित नेटवर्क बीमा फंड $\to$ (4) वैकल्पिक कवरेज दावे में अस्थायी कमी, केवल जब पहले से मौजूद शर्तें और लागू कानून इसे स्पष्ट रूप से अधिकृत करते हैं $\to$ (5) सिद्ध धोखाधड़ी या दुरुपयोग के लिए कानूनी वसूली।

एक कवरेज समायोजन किसी जारीकर्ता के अंतर्निहित वाउचर प्रतिबद्धता को कम नहीं करता है या श्रृंखला पर शेष राशि को नहीं बदलता है, जब तक कि वैध पहले से मौजूद शर्तें और लागू कानून स्पष्ट रूप से उस परिणाम की अनुमति न दें और कोई आवश्यक धारक सहमति प्राप्त न हो।

\textbf{सीमाएं और बहिष्करण।} प्रकाशित कवरेज सीमाओं, पात्र प्रस्तुतियों, साक्ष्य, दावा खिड़कियां, बहिष्कृत मार्ग या घटनाएं, भौगोलिक प्रतिबंध और समाप्त भंडार के उपचार को परिभाषित करेगा। लागू सीमाओं तक पहुंचने के बाद भुगतान शून्य हो सकता है।

\textbf{चित्रणात्मक पुनर्प्राप्ति कार्यक्रम।} यदि अपनाया और प्रकाशित किया गया हो तो:

\begin{enumerate}[leftmargin=*]
\item दावा पहले उत्तरदायी जारीकर्ता या गारंटर बांड से, फिर लागू पूल भंडार से और फिर प्रस्तावित नेटवर्क बीमा कोष से प्राप्त होगा;
\item वैकल्पिक कवरेज का दावा करने के लिए कोई भी कमी उन शर्तों तक सीमित होगी जो पहले से मौजूद कवरेज की शर्तें और लागू कानून, प्रकाशित घटना सीमा तक अधिकृत करते हैं;
\item एक वसूली योजना एक निर्धारित अवधि के लिए वसूले गए मूल्य का एक घोषित हिस्सा लागू कर सकती है, जिसके बाद कोई शेष कवर घाटा सार्वजनिक पोस्ट-मॉर्टम में दर्ज हानि बन जाएगा; और
\item प्रत्येक निर्णय में घटना ID, प्रभावित दावे और वाउचर, निर्णय, वसूली योजना और अपील विंडो के साथ एक रसीद प्रस्तुत की जाएगी।
\end{enumerate}
\subsection*{11.3 गारंटर ढांचा}
\addcontentsline{toc}{subsection}{11.3 गारंटर ढांचा}

इस खंड में जारीकर्ता की जिम्मेदारी, वैकल्पिक पूल सुरक्षा और तृतीय पक्ष गारंटी का अंतर किया गया है।CLC App, CPP, GEF, या किसी भी व्यापक नेटवर्क स्वचालित रूप से एक वाउचर की गारंटी देता है।

\subsection*{बेसलाइन जारीकर्ता की जिम्मेदारी}
\addcontentsline{toc}{subsection}{बेसलाइन जारीकर्ता की जिम्मेदारी}

\begin{itemize}[leftmargin=*]
\item प्रत्येक वाउचर सबसे पहले और मुख्य रूप से इसके जारीकर्ता की जिम्मेदारी है। जारीकर्ता अपने प्रकाशित शर्तों के तहत घोषित वस्तु, सेवा या वैध नकदी समकक्ष प्रदान करने का वादा करता है।
\item जारीकर्ता प्रकाशित करेंगे कि कौन वाउचर प्रस्तुत कर सकता है, वादे की पूर्ति का क्या अर्थ है, यह कहां और कब उपलब्ध है, किस सबूत की आवश्यकता है, और कौन से उपचार लागू होते हैं।
\item यदि कोई जारीकर्ता पूरा नहीं करता है, तो निष्पादक मुख्य जिम्मेदार पक्ष होता है। पूल या नेटवर्क सुरक्षा केवल तब लागू होती है जब इसे अलग से अपनाया जाता है, वित्त पोषित किया जाता है और खुलासा होता है.
\end{itemize}
\subsection*{वैकल्पिक पूल सुरक्षा}
\addcontentsline{toc}{subsection}{वैकल्पिक पूल सुरक्षा}

एक पूल प्रबंधक स्वीकार किए गए वाउचरों के लिए एक संकीर्ण रूप से परिभाषित सुरक्षा जोड़ने का विकल्प चुन सकता है। यह स्वचालित नहीं है और इसे पूल मेटाडेटा और लागू शर्तों में जिम्मेदार पक्ष, वित्त पोषण, पात्र घटनाओं, कैप, खिड़कियों, साक्ष्य, बहिष्करण और उपायों की पहचान करने की आवश्यकता होगी।

उदाहरणात्मक सुरक्षा प्रकारों में निम्नलिखित शामिल हैंः

\begin{enumerate}[leftmargin=*]
\item \textbf{आरक्षित परिसंपत्ति कवरेजः} सत्यापित जारीकर्ता अपर्याप्तता के बाद, जिम्मेदार समूह इकाई एक निर्दिष्ट रिजर्व परिसंपत्ति में एक निश्चित राशि का भुगतान करती है, उसके प्रकाशित कैप और उपलब्ध वित्त पोषित भंडार के अधीन।
\item \textbf{स्विच-बैक विंडोः} एक पात्रता घटना के बाद, पूल पहले या किसी अन्य स्वीकृत परिसंपत्ति में एक समय-सीमित स्वैप पथ प्रदान करता है, जो कैप और इन्वेंट्री के अधीन होता है। यह एक इनवेंटरी-निर्भर तरलता संरक्षण है, न कि एक वादा है कि प्रत्येक स्वॅप reversible है.
\item \textbf{वैकल्पिक वादे की पूर्तिः} उत्तरदायी पक्ष एक प्रकाशित मात्रा या मूल्य सीमा के भीतर एक स्वीकृत प्रतिस्थापन प्रदाता की व्यवस्था करता है।
\item \textbf{विनिमय दर बैंड संरक्षणः} चुनिंदा वाउचर वर्गों के लिए, एक पूल केवल अपने पहले से मौजूद शर्तों में निर्दिष्ट कवरेज समायोजन या स्वैप-बैक उपाय प्रदान करता है। यह जारीकर्ता के अंतर्निहित वाउचर दायित्व को कम नहीं करता है.
\end{enumerate}
\subsection*{सम्भावित वित्तपोषण स्रोत}
\addcontentsline{toc}{subsection}{सम्भावित वित्तपोषण स्रोत}

\begin{itemize}[leftmargin=*]
\item \textbf{जारीकर्ता बांडः} जारीकर्ता द्वारा जमा किया गया प्रतिभूति या खुलासा किए गए रिजर्व में रखा गया और सत्यापित कवर की गई घटना के बाद उपलब्ध है।
\item \textbf{पूल रिजर्वः} जिम्मेदार पूल इकाई द्वारा नियंत्रित संपत्तियां और उसके द्वारा विज्ञापित सुरक्षा के लिए आवंटित।
\item \textbf{थर्ड पार्टी गारंटर बांडः} निर्दिष्ट जारीकर्ताओं, वाउचर वर्गों या घटनाओं के लिए एक ज्ञात बाहरी गारंटर द्वारा रखी गई प्रतिभूति।
\end{itemize}
गारंटर की भागीदारी प्रकाशित पात्रता मानदंडों, बांड आकार, एकाग्रता सीमाओं, निर्णय प्राधिकरण और कानून प्रवर्तन नियमों का पालन करेगी।

\subsection*{दावों की प्रक्रिया}
\addcontentsline{toc}{subsection}{दावों की प्रक्रिया}

स्वीकृत policy audit किए जा सकने वाले triggers बताएगी, जैसे भुनाने की वैध प्रस्तुति के बाद वादा पूरा करने की deadline चूकना, जारीकर्ता का सत्यापित दिवालियापन, covered bridge या escrow failure, या औपचारिक रूप से घोषित incident state।

\begin{itemize}[leftmargin=*]
\item प्रतिभागी दावा कैसे खोलता है और आवश्यक प्रस्तुति और वादे की पूर्ति के प्रमाण प्रदान करता है;
\item जो वाउचर की शर्तों, जारीकर्ता के उत्तरों और तकनीकी रिकॉर्ड को सत्यापित करता है;
\item निर्णय और अपील की खिड़कियां; तथा
\item अधिकृत भुगतान पथ, संपत्ति, कैप और रसीद।
\end{itemize}
जारीकर्ताओं, मध्यस्थता या कानूनी प्रवर्तन से प्राप्त निकासी प्रस्तावित CLC नेटवर्क पूल स्वैप एक्सेस के लिए उपयोग किए जाने से पहले प्रकाशित नीति के अनुसार लागू बांड या भंडार को फिर से भरेंगे।

\subsection*{आवश्यक प्रकटीकरण}
\addcontentsline{toc}{subsection}{आवश्यक प्रकटीकरण}

प्रत्येक कवर किए गए पूल और वाउचर वर्ग के लिए, उत्तरदायी पक्ष प्रकाशित करेगाः

\begin{itemize}[leftmargin=*]
\item क्या एक गारंटर अनुपस्थित है, वैकल्पिक या आवश्यक है;
\item बांड या रिजर्व आकार और एकाग्रता कैप;
\item संरक्षण के प्रकार, संपत्ति, टोपी, खिड़कियां और बहिष्करण;
\item प्रस्तुत करने, पूरा करने, दावा करने और अपील करने की समय सीमा; तथा
\item कौन क्या गारंटी देता है और क्या नहीं।
\end{itemize}
\textbf{संरक्षण सिद्धांत।} पूल प्रबंधकऔर जिम्मेदार कानूनी या शासन संरचनाएं उनके द्वारा विज्ञापित सुरक्षा के लिए जवाबदेह हैं।CPP-संगत तैनाती मानक, रजिस्टर या वैकल्पिक साझा नीतियों प्रदान कर सकते हैं, लेकिन कोई भीCLCऔर न हीGEF स्वचालित रूप से वाउचर या पूल की गारंटी देता है।

\subsection*{11.4 कैप्चर रोधी पहरेदार}
\addcontentsline{toc}{subsection}{11.4 कैप्चर रोधी पहरेदार}

इस टेम्पलेट के अंतर्गत, निम्नलिखित महत्वपूर्ण कार्य होंगे जिन्हें अपनाया गया उच्चतम अनुमोदन स्तर और एक लंबा समय की आवश्यकता होगीः

\begin{enumerate}[leftmargin=*]
\item प्रस्तावित शुल्क जलप्रपात को बदलना, जिसमें इसकी कवरेज और मुख्य संचालन प्राथमिकताएं शामिल हैं;
\item कैनोनिक रजिस्टर की जड़ें बदलना;
\item कवरेज के दायरे, दावों की सीमाओं या निर्णय प्राधिकरण में परिवर्तन;
\item आपातकालीन विराम शक्तियों का विस्तार करना; या
\item इस पत्र में उल्लिखित फोर्कलिटी, पारदर्शिता या पूल संप्रभुता प्रतिबद्धताओं को कमजोर करना।
\end{enumerate}
\subsection*{11.5 फोर्क और बाहर निकलने की प्रक्रिया}
\addcontentsline{toc}{subsection}{11.5 फोर्क और बाहर निकलने की प्रक्रिया}

यदि शासन को कैप्चर किया गया या मूल्यों ने महत्वपूर्ण रूप से बहाली की, समुदायों, पूल प्रबंधक और ऑपरेटर नेटवर्क प्रशासन परत को फोर्क करके बाहर निकलने का प्रयास कर सकते हैं। अंतर्निहित पूल और वाउचर की निरंतरता तैनात अनुबंधों, कुंजीओं, इंटरफेस, बुनियादी ढांचे, तृतीय-पक्ष सेवाओं और लागू दायित्वों पर निर्भर करेगी।

बाहर निकलने की प्रक्रियाः

\begin{enumerate}[leftmargin=*]
\item \textbf{एक स्नैपशॉट प्रकाशित करेंः} चयनित रजिस्टर, वाउचर, मूल्यों, सीमाओं और शुल्क नीतियों का निर्यात करें, फिर एक हस्ताक्षरित स्नैपशॉट हैश प्रकाशित करें।
\item \textbf{शासन सेवाओं को पुनः तैनात करनाः} नए रजिस्ट्री रूट, मार्ग सेवाएं और किसी भी स्वीकृत शुल्क या कवरेज मॉड्यूल को एक नई उत्तरदायी संरचना के तहत तैनात करना।
\item \textbf{पुनः पंजीकरणः} पूल प्रबंधक को नए रूट के तहत अपने पूल पते को पंजीकृत करके अपना विकल्प चुनने की अनुमति दें, बिना किसी अन्य प्रकार के कार्यात्मक वाउचर को स्थानांतरित करने के लिए धारकों की आवश्यकता के।
\item \textbf{पुनः नियुक्त करने वाले ग्राहकः} एसडीके और इंटरफेस में एक चयन योग्य नेटवर्क प्रोफ़ाइल के रूप में नई रूट जोड़ें, खुलासा किए गए शासन प्रक्रिया के माध्यम से किसी भी डिफ़ॉल्ट परिवर्तन के साथ।
\item \textbf{एक पुल अवधि का प्रबंधन करेंः} सुरक्षित स्थानों पर संगत मार्ग बनाए रखें और नए प्रोफाइल के नियमों का उल्लंघन करने वाले मार्गों को अस्वीकार करें।
\end{enumerate}
डिजाइन उद्देश्य यह है कि एक कैनोनिक रजिस्ट्री छोड़ने से अन्यथा कार्यात्मक स्थानीय पूल निष्क्रिय नहीं होते हैं। वास्तविक निरंतरता तैनाती पर निर्भर रहती है; संघ एक ऑप्ट-इन डिस्कवरी और समन्वय परत है।


\section*{12. प्रस्तावित तरलता कार्यक्रम अवधि पत्र (गैर-बाध्यकारी रूपरेखा)}
\addcontentsline{toc}{section}{12. प्रस्तावित तरलता कार्यक्रम अवधि पत्र (गैर-बाध्यकारी रूपरेखा)}

यह भविष्य के लिए एक स्पष्ट चेकलिस्ट है, अलग से प्रलेखित कार्यक्रम। यह एक प्रस्ताव नहीं है, एक वर्तमान ऐप सुविधा, या तरलता का वादा, पहुंच, बीमा, रिटर्न, मूल्य, या बाहर निकलना।

\begin{itemize}[leftmargin=*]
\item \textbf{पात्र योगदान:} कार्यक्रम के अनुसार स्थिर सिक्के, रिजर्व टोकन या अनुमोदित सामुदायिक वाउचर।
\item \textbf{स्थान:} नामितप्रतिबद्धता पूलया प्रस्तावितCLCएक अलग से अपनाए गए जनादेश के तहत नेटवर्क पूल।
\item \textbf{तालाबंदी और बाहर निकलनाः} कार्यक्रम विशिष्ट शर्तें, जैसे रोलिंग 30--90 दिन की अवधि और तनाव के तहत खुलासे किए गए द्वार।
\item \textbf{नीतिगत शुल्क क्रेडिटः} प्रकाशित कैप और विंडो के तहत एक वैकल्पिक एक्स-पोस्ट तंत्र, न कि वादा किया गया रिटर्न।
\item \textbf{जोखिम साझा करनाः} वैकल्पिक कवरेज का दावा करने के लिए किसी भी कमी को पहले से मौजूद शर्तों और लागू कानून द्वारा स्पष्ट रूप से अनुमोदित किया जाना चाहिए। यह स्वयं किसी जारीकर्ता की वाउचर प्रतिबद्धता या श्रृंखला पर शेष राशि को नहीं बदलता है।
\item \textbf{रिपोर्टिंगः} नियमित डैशबोर्ड और प्रमाणपत्र जो पूल के स्वास्थ्य, इन्वेंट्री, सीमाओं, शुल्क और किसी भी वित्त पोषित कवरेज रिजर्व को कवर करते हैं।
\item \textbf{करारः} आय के उपयोग पर प्रतिबंध, अध्यात्म नियंत्रण, सीमा स्तर, संघर्ष नियम और उपाय।
\item \textbf{प्रस्तावित टोकन की संभावित पात्रताः} उन योगदानकर्ताओं के लिए पात्र हो सकते हैं जिनके पास पूल नामित पूल है, जो मापे गए पूल-स्वैप एक्सेस और जारीकर्ता की पुष्टि पूर्णता पर आधारित प्रस्तावित CLC गवर्नेंस टोकन अनुदानों के लिए अर्हता प्राप्त कर सकते हैं, स्वीकार किए गए शर्तों, एंटी-गेमिंग नियमों और नीतिगत सीमाओं के अधीन।
\end{itemize}


\section*{13. सरल भाषा का शब्दकोश}
\addcontentsline{toc}{section}{13. सरल भाषा का शब्दकोश}

\begin{itemize}[leftmargin=*]
\item \textbf{Cosmo-Local Credit (CLC):} CLC App सहित उत्पाद परिवार, प्रलेखन और व्यापक प्रतिबद्धता साझाकरण कार्य।
\item \textbf{CLC App:} प्रगतिशील वेब ऐप द्वारा संचालित GEF पर \texttt{cosmolocal.credit}.
\item \textbf{प्रतिबद्धता संवर्धन प्रोटोकॉल (CPP):} रखरखाव, मूल्यांकन, सीमांकन, विनिमय और उत्तरदायी शासन का पुनः प्रयोज्य मॉडल।
\item \textbf{Protocol v1.1.0:} सत्यापित सार्वजनिक स्मार्ट अनुबंध रिलीज.
\item \textbf{प्रतिबद्धता पूल:} एक नियंत्रित व्यवस्था जो समर्थित परिसंपत्तियों को स्वीकार करती है और विनिमय नियम प्रकाशित करती है।
\item \textbf{\texttt{SwapPool}:} वर्तमान टोकन वॉल्ट और प्रत्यक्ष स्वैप अनुबंध।
\item \textbf{टोकन:} केवल token mechanics अपने आप भुनाए जा सकने वाला वादा नहीं बनाती।
\item \textbf{वाउचर:} ऐसा token या record जिसे जारीकर्ता प्रकाशित शर्तों के तहत भुनाए जा सकने वाले वादे के रूप में प्रस्तुत करता है।
\item \textbf{प्रस्ताव:} वाउचर से जुड़ी वस्तुओं या सेवाओं के लिए कैटलॉग रिकॉर्ड।
\item \textbf{जारीकर्ता:} वाउचर प्रतिबद्धता को प्रकाशित करने और पूरा करने के लिए जिम्मेदार पक्ष।
\item \textbf{पूल प्रबंधक:} पूल के नियमों को प्रकाशित और प्रशासित करने वाली जिम्मेदार संरचना।
\item \textbf{पूल विनिमय दर या बोलीः} एक कॉन्फ़िगर किए गए कोटर द्वारा उत्पादित लेनदेन पैरामीटर; नकद या भुनाने मूल्य का प्रमाण नहीं।
\item \textbf{पूल टोकन बैलेंस कैपः} एक पूल द्वारा रखे गए टोकन के लिए वर्तमान \texttt{Limiter} अधिकतम। ऐप इसे लेबल कर सकता है \textbf{``क्रेडिट सीमा।''}
\item \textbf{पूल स्वैप:} एक \texttt{SwapPool} के माध्यम से प्रत्यक्ष परिसंपत्ति विनिमय।
\item \textbf{स्वैप सेटलमेंटः} पूल स्वैप को पूरा करने वाले ऑन-चेन हस्तांतरण और शुल्क लेखांकन।
\item \textbf{मुआवजे की प्रस्तुतिः} जारीकर्ता को वाउचर इकाइयों की वापसी या अन्यथा प्रस्तुत करना।
\item \textbf{वादे की पूर्ति:} जारीकर्ता वादा किया गया वस्तु, सेवा, लाभ या अन्य प्रदर्शन प्रदान करता है।
\item \textbf{दायित्व की समाप्तिः} पूर्ण इकाइयों को फिर से प्रस्तुत करने से रोकने के लिए एक रिकॉर्ड।
\item \textbf{प्रस्तावित निष्पादन राउटरः} भविष्य के सॉफ्टवेयर जो मार्गों को अधिकृत और निष्पादित करेगा। वर्तमान \texttt{SwapRouter} केवल उद्धरण करता है।
\item \textbf{प्रस्तावित CLC नेटवर्क पूलः} नेटवर्क स्तर पर भविष्य की क्लियरिंग और बजट समन्वय व्यवस्था।
\item \textbf{प्रस्तावित CLC शासन टोकनः} भविष्य के शासन डिजाइन, न कि वर्तमान ऐप या Protocol v1.1.0 संपत्ति।
\item \textbf{नेटवर्क रैकः} भविष्य के नेटवर्क बजट में स्थानांतरित किए जाने वाले पूल शुल्क का प्रस्तावित हिस्सा।
\item \textbf{गारंटी या बीमाः} एक ऐसा संरक्षण जो केवल तभी लागू होता है जब किसी पहचाने गए पक्ष ने इसे स्पष्ट रूप से ग्रहण किया हो, वित्तपोषित किया हो और प्रकाशित किया हो।
\end{itemize}

\section*{14. प्रस्तावित KPIs और स्वास्थ्य संकेतक}
\addcontentsline{toc}{section}{14. प्रस्तावित KPIs और स्वास्थ्य संकेतक}

एक तैनाती को केवल तब KPI प्रकाशित करना चाहिए जब वह घटना, स्रोत, समूह या अवधि, इकाई, मूल्यांकन विधि और समयशीर्षक, बहिष्करण, सुधार नीति, आउट-चेन साक्ष्य और डेटा गुणवत्ता सीमाओं को परिभाषित कर सके।

प्रस्तावित उपायों में निम्नलिखित शामिल हैंः

\begin{itemize}[leftmargin=*]
\item वैध प्रतिपूर्ति प्रस्तुतियाँ;
\item कोहर्ट आधारित वादे की पूर्ति दर;
\item डिस्चार्ज की पूर्णता;
\item प्रस्तुति से वादे की पूर्ति तक की विलंबता;
\item अधिग्रहण से प्रस्तुति तक holding की अवधि;
\item एक निर्धारित पद्धति के तहत अपव्यय योग्य प्रतिबद्धताएं;
\item पूल स्वैप की पूर्ण मात्रा;
\item मापा गया पूल इन्वेंट्री;
\item पूल टोकन बैलेंस कैप का उपयोग;
\item उद्धरण पास दर, मार्ग निष्पादन दर से अलग रखी जाती है;
\item स्पष्ट रूप से कवर किए गए जोखिम द्वारा विभाजित पात्र भंडार;
\item पहचान की गई पॉलिसी के तहत गारंटर का दावा और वसूली;
\item प्रस्तावित नेटवर्क रैक और सेवा शुल्क वास्तव में प्राप्त किए गए, डबल-कंटिंग सकल पूल शुल्क के बिना;
\item शासन का पता लगाना, निर्णय लेना, विराम देना, मरम्मत करना और बंद करना;
\item शिकायतें, विवाद, सुधार और उपचार; तथा
\item अलग-अलग सामाजिक या पारिस्थितिक परिणामों को दर्शाया।
\end{itemize}
ऑन-चेन लेनदेन डेटा स्वयं पहचान, जारीकर्ता प्रदर्शन, संतुष्टि, कारण संबंधीता, सामुदायिक स्वास्थ्य या सामाजिक प्रभाव को साबित नहीं करता है। इन दावों के लिए अपनी पद्धति और सबूत की आवश्यकता होती है।

देखिये \href{https://docs.cosmolocal.credit/hi/white-paper/appendix-c-kpi-definitions}{परिशिष्ट सी} प्रस्तावित माप विनिर्देश के लिए।


\section*{15. रोडमैप (सूचक)}
\addcontentsline{toc}{section}{15. रोडमैप (सूचक)}

\subsection*{15.1 वर्तमान आधार}
\addcontentsline{toc}{subsection}{15.1 वर्तमान आधार}

GEF संचालित करता हैCLC App पर \texttt{cosmolocal.credit} Gnosis Chain पर ऐप समर्थित खाते और वॉलेट एक्सेस, सार्वजनिक बाजार कैटलॉग, और टोकन, वाउचर, ऑफर के लिए इंटरफेस प्रदान करता है।प्रतिबद्धता पूल, हस्तांतरण, और प्रत्यक्ष पूल स्वैप।

Protocol v1.1.0 वर्तमान अनुबंध की आधारशिला प्रदान करता हैः\texttt{GiftableToken}, प्रत्यक्ष \texttt{SwapPool} निष्पादन, वैकल्पिक रजिस्टर, मूल्यांकन मॉड्यूल, पूल टोकन-बालेंस कैप, पूल शुल्क, अतिरिक्त प्रोटोकॉल शुल्क और केवल उद्धरण \texttt{SwapRouter}. उपलब्धता अभी भी इंटरफ़ेस, तैनात अनुबंधों, सूची, विन्यास, नेटवर्क की स्थिति, पात्रता, प्रदाताओं, अधिकार क्षेत्र और प्रकाशित जारीकर्ता या पूल शर्तों पर निर्भर करती है।

\subsection*{15.2 प्रस्तावित मील के पत्थर}
\addcontentsline{toc}{subsection}{15.2 प्रस्तावित मील के पत्थर}

ये नामित मील के पत्थर डिजाइन निर्देश हैं, रिलीज नंबर, डिलीवरी की तारीखें या प्रतिबद्धताएं नहीं जो एक सुविधा लॉन्च करेगी।

\begin{itemize}[leftmargin=*]
\item \textbf{शासन फाउंडेशन} प्रस्तावितCLCशासन टोकन; प्रस्तावितCLCनेटवर्क पूल; शुल्क एडाप्टर; क्वारोम, टाइमलॉक और शासन नींव।
\item \textbf{रूटिंग और अवलोकनशीलता:} राउटर SDK और रजिस्ट्री एपीआई; स्वास्थ्य डैशबोर्ड; एक प्रस्तावित बीमा नीति ढांचा; ऑप्ट-इन rebalancing इरादे; बैच नेट प्रोटोटाइप।
\item \textbf{Cross-Domain Risk Tools:} HTLC या एस्क्रो रूटिंग; अलग से नियंत्रित गारंटर मॉड्यूल; रोलिंग, खाता और स्तरबद्ध सीमा पूर्वनिर्धारित।
\item \textbf{विनियमित पहुंचः} तैनाती-विशिष्ट तृतीय-पक्ष भुगतान सेवाएं; व्यक्तिगत सूक्ष्म पूल; अनुपालन-सेवा का पता लगाना; वाउचर वर्गों के तीसरे पक्ष के ऑडिट।
\end{itemize}
किसी भी भुगतान सेवा को लागू अधिकार क्षेत्र, पात्रता, शुल्क, सीमाएं और शर्तों के तहत अलग से पहचाने गए तृतीय पक्ष द्वारा प्रदान किया जाएगा।CLC App और Protocol v1.1.0 अनुबंध स्वयं फिएट रेल का संचालन नहीं करते हैं।

बहु-हॉप निष्पादन, साझा बीमा, प्रस्तावितCLCशासन का प्रतीक और प्रस्तावितCLCनेटवर्क पूल, बैच नेटिंग, व्यक्तिगत माइक्रो-पूल और विनियमित भुगतान सेवाएं तब तक प्रस्तावित या तैनाती से निर्भर रहती हैं जब तक कि एक कार्यान्वयन और उसके शासन शर्तें उन्हें सक्रिय नहीं पहचानती हैं।


\section*{16. प्रस्तावित मूल्यांकन टेम्पलेट}
\addcontentsline{toc}{section}{16. प्रस्तावित मूल्यांकन टेम्पलेट}

एक रजिस्ट्री, पूल या भविष्य की तरलता कार्यक्रम इस टेम्पलेट को अनुकूलित कर सकता है। यह वर्तमान सार्वभौमिक लिस्टिंग मानक या गारंटी नहीं है।

\begin{enumerate}[leftmargin=*]
\item \textbf{देखे गए तथ्यः} जारीकर्ता की पहचान और इतिहास, वाउचर शर्तें, ऑफ़र, क्षमता प्रमाण, पूल कॉन्फ़िगरेशन, नियंत्रक, इन्वेंट्री, सीमाएं, भंडार, गारंटर, घटनाएं और अनिर्णित दायित्व।
\item \textbf{उद्देश्य और प्रभावित पक्षः} अपेक्षित लाभ, जोखिम, संघर्ष, पारिस्थितिक और सामाजिक प्रभाव और कौन नुकसान या जिम्मेदारी उठाता है।
\item \textbf{मूल्यांकनः} प्रस्ताव मूल्य, वाउचर घोषित मूल्य, पूल विनिमय दर विधि, ओरेकल संदर्भ, इकाइयां, समय टिकट और किसी भी विचलन के कारण।
\item \textbf{सीमाएँः} पात्रता, प्रतिबंधित उपयोग, एकाग्रता, टोकन बैलेंस कैप, विराम, भौगोलिक या कानूनी प्रतिबंध और समीक्षा ट्रिगर।
\item \textbf{संरक्षण:} निर्दिष्ट बाध्य पक्ष, कवर की गई घटना, वित्तपोषण, सीमा, बहिष्करण, अवधि, साक्ष्य, दावा प्रक्रिया और उपचार।
\item \textbf{निर्णयः} मानव समीक्षा के लिए अनुमोदन, संशोधन, अस्वीकार, विराम या वृद्धि; निर्णय लेने वाले, कारणों, शर्तों, समीक्षा तिथि और अपील या सुधार प्रक्रिया की पहचान करें।
\end{enumerate}
आउटपुट में क्युरेशन, लिमिट, रिजर्व या सत्यापन को अनुमोदन, जारीकर्ता क्षमता, बीमा या गारंटीकृत वादे की पूर्ति के रूप में वर्णित नहीं किया जाना चाहिए जब तक कि जिम्मेदार पक्ष ने स्पष्ट रूप से उस संरक्षण की पुष्टि नहीं की है।


\section*{17. कानूनी और अनुपालन नोट}
\addcontentsline{toc}{section}{17. कानूनी और अनुपालन नोट}

CPP टोकन, वाउचर, पूल और स्वैप का समन्वय कर सकता है जिनका कानूनी उपचार उनके डिजाइन, विपणन, उपयोग, जिम्मेदार पक्षों और अधिकार क्षेत्र पर निर्भर करता है। इस श्वेत पत्र में कुछ भी कानूनी, कर, निवेश, क्रेडिट या वित्तीय सलाह नहीं है।

CLC App का प्रयोग \href{https://docs.cosmolocal.credit/hi/governance/terms}{सेवा की शर्तें}. GEF ऐप और सहायक बुनियादी ढांचे का संचालन करता है।GEF स्पष्ट रूप से एक और भूमिका निभाता है, यह issuer नहीं है,पूल प्रबंधक, संरक्षक, ऋणदाता, उधारकर्ता, दलाल, गारंटर, रिडीमर, सलाहकार, बीमाकर्ता या प्रतिभागियों के बीच दायित्वों का पक्ष।

तैनाती-निर्भर भुगतान सेवाओं को अपने जिम्मेदार प्रदाता, अधिकार क्षेत्र, पात्रता, शुल्क, सीमाएं, रखरखाव मॉडल और शर्तों की पहचान करनी चाहिए। एक कनेक्टर की उपस्थिति का मतलब यह नहीं है कि GEF या एक पूल अंतर्निहित विनियमित सेवा संचालित करता है।

\subsection*{17.1 वाउचर और ऑफर का खुलासा}
\addcontentsline{toc}{subsection}{17.1 वाउचर और ऑफर का खुलासा}

एक जारीकर्ता को प्रकाशित करना चाहिए और अद्यतित रखना चाहिएः

\begin{itemize}[leftmargin=*]
\item जिम्मेदार पहचान और संपर्क जानकारी;
\item संबंधित प्रस्ताव, इकाई, आपूर्ति, घोषित मूल्य और क्षमता;
\item प्रस्तुति के स्थान और प्रक्रियाएं;
\item वादे की पूर्ति का समय और प्रमाण;
\item समाप्ति, शुल्क, कर, प्रतिबंध और सामग्री परिवर्तन के नियम;
\item शिकायतें, प्रतिस्थापन, धनवापसी या अन्य उपचार; तथा
\item पूरा होने के बाद पुनः उपयोग को रोकने वाली डिस्चार्ज विधि।
\end{itemize}
एक टोकन अनुबंध इन शर्तों को स्थापित नहीं करता है या प्रदर्शन का प्रमाण नहीं देता है।

\subsection*{17.2 पूल प्रकटीकरण}
\addcontentsline{toc}{subsection}{17.2 पूल प्रकटीकरण}

एक पूल प्रबंधक को प्रकाशित करना चाहिएः

\begin{itemize}[leftmargin=*]
\item पूल का उद्देश्य और जिम्मेदार निर्णय लेने वाले;
\item \texttt{SwapPool} के मालिक, प्रॉक्सी प्रशासक, निर्भरता नियंत्रक और शुल्क प्राप्तकर्ता;
\item अनुमोदित संपत्ति और निलंबन या हटाने के नियम;
\item मूल्य निर्धारण विधियां, उद्धरण स्रोत, शुल्क, टोकन बैलेंस कैप, इन्वेंट्री और मालिक निकासी शक्तियाँ;
\item योगदान और बहिष्कार अधिकार;
\item प्रत्येक आरक्षित, गारंटी, गारंटर, बीमा पॉलिसी और हानि आवंटन नियम; तथा
\item शासन, उन्नयन, आपातकाल, शिकायत, प्रवासन और समाप्ति प्रक्रियाएं।
\end{itemize}
सूचीकरण GEF द्वारा अनुमोदन, मूल्यांकन, बीमा या गारंटी नहीं है।

\subsection*{17.3 कार्रवाई और दायित्व}
\addcontentsline{toc}{subsection}{17.3 कार्रवाई और दायित्व}

साधारण \textbf{पूल स्वैप} दिखाई गई quote और transaction limits के तहत समर्थित संपत्तियों का आदान-प्रदान करता है। संपत्ति की दिशा उसे ऋण, repayment, जारीकर्ता द्वारा वाउचर भुनाना या वास्तविक दुनिया में वादा पूरा करना नहीं बनाती।

\textbf{मुआवजे की प्रस्तुति} किसी जारीकर्ता को वाउचर इकाइयां लौटाता है या प्रस्तुत करता है। \textbf{ पूरा होना} जारीकर्ता का वादा किया गया प्रदर्शन है। \textbf{ डिस्चार्ज } एप्लिकेशन की पूर्ण इकाइयों के रिकॉर्ड ताकि वे पुनः उपयोग नहीं किया जा सकता है। \textbf{`भुनाएँ'} कार्रवाई वर्तमान में टोकन के मालिक को हस्तांतरण की तैयारी करती है; यह स्थानांतरण अकेले वादे की पूर्ति साबित नहीं करता है।

ऐप की \textbf{`रिटायर वाउचर'} कार्रवाई सूचीबद्ध निष्कासन है। यह शेष राशि को नहीं जलाता, दावों को रद्द करता है, या दायित्वों का दायित्व की समाप्ति नहीं करता है।

भविष्य के ऋण या अन्य क्रेडिट उत्पाद को उपयोग से पहले अलग-अलग प्रस्तुत किए गए पूरक और लेनदेन शर्तों की आवश्यकता होगी। कोई सामान्य भेज, योगदान, पूल जमा, पूल स्वैप, प्रस्तुति, वादे की पूर्ति या डिस्चार्ज केवल इसके दिशा या परिसंपत्ति प्रकार के कारण ऋण नहीं बनाता है।

\subsection*{17.4 सुरक्षा, साक्ष्य और स्थानीय कानून}
\addcontentsline{toc}{subsection}{17.4 सुरक्षा, साक्ष्य और स्थानीय कानून}

एक सीमा, रिजर्व, रजिस्ट्री प्रविष्टि, ऐप लिस्टिंग या ब्लॉकचेन लेनदेन स्वचालित रूप से गारंटी या बीमा पॉलिसी नहीं है। किसी भी सुरक्षा को इसके जिम्मेदार पक्ष, कवर किए गए घटना, वित्त पोषण, कैप, बहिष्करण, अवधि, सबूत, दावा प्रक्रिया और लागू शर्तों की पहचान करनी चाहिए।

सार्वजनिक श्रृंखला साक्ष्य पते, परिसंपत्तियों, राशि, समय टिकट और अनुबंध घटनाओं को दिखा सकते हैं। यह स्वयं पहचान, जारीकर्ता क्षमता, वादे की पूर्ति, संतुष्टि, शासन विचार-विमर्श, कानूनी डिस्चार्ज या सामाजिक प्रभाव साबित नहीं करता है।

कानून, प्रतिबंधों, पात्रता, अनुबंध राज्य, इन्वेंट्री, सीमाओं, सुरक्षा, अधिकार क्षेत्र या तृतीय-पक्ष सेवाओं के कारण सुविधाएं प्रतिबंधित या अनुपलब्ध हो सकती हैं। जारीकर्ता, पूल प्रबंधक, प्रदाता और प्रतिभागियों को लागू कानून का निर्धारण करने और उनका पालन करने की जिम्मेदारी बनी रहती है।


\section*{18. निष्कर्ष}
\addcontentsline{toc}{section}{18. निष्कर्ष}

Cosmo-Local Credit एक वर्तमान आवेदन को जोड़ता है और Protocol v1.1.0 स्वतंत्र रूप से शासित के लिए एक व्यापक प्रस्तावित डिजाइन के साथ नींवप्रतिबद्धता पूल. वर्तमान प्रत्यक्ष पूल स्वैप, कोट और लिमिट घटक, तथा सार्वजनिक लेनदेन रिकॉर्ड उत्तरदायी आदान-प्रदान का समर्थन कर सकते हैं, लेकिन वे जारीकर्ता की वादे की पूर्ति साबित नहीं करते, योगदानकर्ता अधिकार पैदा नहीं करते या मूल्य, तरलता, मुआवजा, बीमा, कानूनी स्थिति या हानि से सुरक्षा की गारंटी नहीं देते।

प्रस्तावित निष्पादन रूटिंग, नेटिंग, गवर्नेंस एसेट, नेटवर्क क्लीरिंग, तरलता कार्यक्रम और इस पेपर में वर्णित साझा सुरक्षा के लिए अलग कार्यान्वयन, वित्तपोषण, शासन और प्रकाशित शर्तों की आवश्यकता होगी। डिजाइन का उद्देश्य जारीकर्ता प्रदर्शन, पूल प्रशासन और पहचाने गए पक्षों के साथ वास्तविक दुनिया की जवाबदेही को बनाए रखते हुए सहकार्यशीलता का विस्तार करना है।


\section*{ए. प्रस्तावित माप और शुल्क मॉडल}
\addcontentsline{toc}{section}{ए. प्रस्तावित माप और शुल्क मॉडल}

इस परिशिष्ट में एक प्रस्तावित माप ढांचे को परिभाषित किया गया है। Protocol v1.1.0 जारीकर्ता वादे की पूर्ति, वास्तविक दुनिया के डिस्चार्ज या नीचे आवश्यक प्रत्येक डेटा फ़ील्ड का रिकॉर्ड नहीं करता है।

\subsection*{घटना और स्टॉक की परिभाषा}
\addcontentsline{toc}{subsection}{घटना और स्टॉक की परिभाषा}

वाउचर वर्ग \emph{j}, कोहोर्ट या अवधि \emph{t} के लिए, और खुलासा मूल्य निर्धारण विधि \emph{m}:

\begin{itemize}[leftmargin=*]
\item \texttt{O\_\{j,t,m\}}: माप सीमा पर बकाया पात्र प्रतिबद्धताओं का मूल्य।
\item \texttt{X\_\{j,t,m\}}: अवधि के दौरान संपन्न पूल स्वैप का मूल्य।
\item \texttt{P\_\{j,t\}}: निकासी के लिए जारीकर्ता को मान्य रूप से प्रस्तुत किए गए इकाइयां।
\item \texttt{F\_\{j,t\}}: अलग से प्रमाणित जारीकर्ता वादे की पूर्ति वाली प्रस्तुत इकाइयां।
\item \texttt{G\_\{j,t\}}: पूर्ण इकाइयां जिनके डिस्चार्ज रिकॉर्ड के साथ पुनः उपयोग की रोकथाम होती है।
\end{itemize}
\texttt{O} को केवल टोकन की आपूर्ति से नहीं निकाला जा सकता है। एक माप नीति में जिम्मेदार जारीकर्ता का पता लगाना चाहिए और लागू होने पर, emitter-holding inventory, जलाए गए इकाइयों, expired units, discharged units, test balances, inaccessible balances और tokens को बाहर करना चाहिए जिनकी शर्तें किसी third party के लिए अप्रासंगिक प्रतिबद्धता पैदा नहीं करती हैं.

प्रत्येक मूल्यवान उपाय को भिन्न पूल विनिमय दरों की इकाई, स्रोत, मूल्य निर्धारण विधि, समयशीर्षक और उपचार प्रकाशित करना चाहिए।\texttt{X} या प्रस्तुत सबूत; यह स्वयं नहीं स्थापित करता है\texttt{F} या\texttt{G}.

\subsection*{कोहर्ट आधारित वादे की पूर्ति उपाय}
\addcontentsline{toc}{subsection}{कोहर्ट आधारित वादे की पूर्ति उपाय}

वैध मुआवजा प्रस्तुतियों के एक समूह के लिएः

\texttt{FulfillmentRate = FulfilledPresentments / ValidPresentments}

\texttt{DischargeCompleteness = DischargedFulfillments / FulfilledPresentments}

प्रत्येक अंकन और संकेतक में एक ही बंद या परिपक्व कोहोर्ट का उपयोग करें। रिपोर्ट अस्वीकृत, वापस ले लिया गया, समाप्त हो गया, विवादित, आंशिक रूप से पूरा हुआ, सही किया गया और अभी भी खुला प्रस्तुति अलग-अलग।

\texttt{FulfillmentLatency = median(t\_fulfilled - t\_presented)}

\texttt{HoldingDuration = median(t\_presented - t\_acquired)}

निष्पादन विलंबता जारीकर्ता सेवा को प्रस्तुत करने के बाद मापती है। रखरखाव अवधि एक अलग उपाय है और इसे रिडीम लटेंसी के रूप में लेबल नहीं किया जाना चाहिए।

\subsection*{वेग के भिन्न उपाय}
\addcontentsline{toc}{subsection}{वेग के भिन्न उपाय}

एक प्रस्तावित प्रतिबद्धता-निष्कासन गति की गणना केवल तभी की जा सकती है जब \texttt{O} और पूर्ति मूल्य का उपयोग उसी मूल्यांकन पद्धति से किया जाता हैः

\[
V_{\text{\latinfont commit},t} = \frac{\text{\latinfont FulfilledValue}_t}{\text{\latinfont AverageOutstandingEligibleCommitmentValue}_t}
\]

एक पूल स्वैप गतिविधि उपाय अलग हैः

\texttt{V\_swap,t = PoolSwapValue\_t / AverageMeasuredPoolInventoryValue\_t}

कोई भी मूल्य सामाजिक प्रभाव, जारीकर्ता क्षमता, लाभप्रदता या नकदी परिवर्तनीयता का प्रमाण नहीं है।

\subsection*{प्रस्तावित नेटवर्क राजस्व}
\addcontentsline{toc}{subsection}{प्रस्तावित नेटवर्क राजस्व}

लेटः

\begin{itemize}[leftmargin=*]
\item \texttt{PF\_t} अवधि के दौरान उत्पन्न सकल पूल शुल्क है;
\item \texttt{NR\_t} प्रस्तावित नेटवर्क रैक हो जो वास्तव में उन पूल शुल्क के खुलासा किए गए हिस्से के रूप में प्राप्त हुआ है;
\item \texttt{RF\_t} अलग से प्रस्तावित रूटिंग या सेवा शुल्क वास्तव में प्राप्त; तथा
\item \texttt{$\chi$\_t} प्राप्त राजस्व का मापा गया हिस्सा है जो लागतों और नीतिगत बाधाओं के बाद घोषित नकदी में नामित उपयोग के लिए पात्र और परिवर्तनीय है।
\end{itemize}
प्रस्तावित नेटवर्क बजट परिप्रेक्ष्य सेः

\texttt{ProposedNetworkRevenue\_t = NR\_t + RF\_t}

\texttt{CashUsableNetworkRevenue\_t = $\chi$\_t $\times$ (NR\_t + RF\_t)}

जोड़ना नहीं\texttt{PF\_t} करने के लिए \texttt{NR\_t}: रैक सकल पूल शुल्क से हस्तांतरण है और अन्यथा दो बार गिना जाएगा।Protocol v1.1.0 इसके बजाय एक अतिरिक्त प्रोटोकॉल शुल्क का समर्थन करता है; इसकी प्राप्तियां इस प्रस्तावित रैक मॉडल से अलग रिपोर्ट की जानी चाहिए।


\section*{B. भविष्य शुल्क जलप्रपात}
\addcontentsline{toc}{section}{B. भविष्य शुल्क जलप्रपात}

यह भविष्य की तैनाती के लिए एक प्रस्तावित डिजाइन है। यह केवल तभी लागू होता है जब इसे लागू, वित्त पोषित और प्रकाशित शासन और प्रतिभागी शर्तों के माध्यम से अपनाया जाता है। इसमें वर्तमान Protocol v1.1.0 शुल्क, योगदानकर्ता अधिकार, रिजर्व, गारंटी या बीमा व्यवस्था का वर्णन नहीं किया गया है।

\texttt{F\_in} को एक युग के दौरान प्रस्तावित नेटवर्क बजट द्वारा वास्तव में प्राप्त राजस्व माना जाएः प्रस्तावित नेट रैक, अलग रूटिंग या सेवा शुल्क और कोई अन्य स्पष्ट रूप से निर्दिष्ट राजस्व।

राजस्व परिसंपत्तियों को निम्नानुसार वर्गीकृत किया जाना चाहिएः

\begin{itemize}[leftmargin=*]
\item \texttt{E\_cash}: स्वीकार की गई नीति के तहत घोषित नकदी मुद्रा में उपयोग के लिए पात्र संपत्ति; या
\item \texttt{E\_kind}: नकद रूप में परिवर्तनीय नहीं होने वाले वस्तुगत वाउचर या अन्य परिसंपत्तियां।
\end{itemize}
किसी भी रूपांतरण नीति के लिए परिसंपत्तियों और स्थान अनुमतियों, जिम्मेदार अधिकारियों, मूल्य स्रोतों, स्लिप-अप सीमाओं, रिपोर्टिंग और लागू कानूनी नियंत्रण की आवश्यकता होगी।

एक प्रस्तावित झरना इस क्रम में प्राप्त राजस्व लागू कर सकता हैः

\begin{enumerate}[leftmargin=*]
\item \textbf{कवर किए गए रिजर्व लक्ष्यः} केवल परिभाषित कवर किए गए जोखिम, पात्र परिसंपत्तियों, बहिष्करणों और दावा नियमों के आधार पर अपनाया गया रिजर्व या बीमा लक्ष्य निधि।
\item \textbf{कोर ऑपरेशनः} एक खुलासा किया गया, सीमित परिचालन बजट निधि।
\item \textbf{लिक्विडिटी जनादेशः} अनुमोदित, अलग से नियंत्रित पूल या रूटिंग कार्यक्रम।
\item \textbf{शेष बजट:} प्रकाशित सीमाओं के तहत संचालन, तरलता कार्यक्रमों और एक अतिरिक्त बफर के बीच शेष राशि आवंटित करें।
\end{enumerate}
आरक्षित लक्ष्य स्वयं गारंटी नहीं है। किसी भी बीमा या गारंटी को बाध्य पक्ष, कवर की गई घटना, वित्तपोषण, सीमा, बहिष्करण, अवधि, साक्ष्य, दावा प्रक्रिया और हानि आवंटन का पता लगाना चाहिए।

\textbf{गार्डरेल:} जलप्रपात आवंटन स्वीकृत कवरेज, संचालन और तरलता सेवाओं के लिए है। इन्हें मूल्य समर्थन कार्यों के रूप में ढांचाबद्ध या निष्पादित नहीं किया जाना चाहिए।

प्रस्तावित मॉडल के तहत, एक अलग से सक्षम बाहरी तरलता कार्यक्रम के माध्यम से प्राप्त किसी भी प्रस्तावित CLC शासन टोकन को रिटायर किया जाएगा या एक खुलासा किए गए गैर-मतदान सिंक में रखा जाएगा। यह तंत्र तैनात नहीं किया गया है।


\section*{C. प्रस्तावित KPI परिभाषाएं}
\addcontentsline{toc}{section}{C. प्रस्तावित KPI परिभाषाएं}

ये KPIs एक प्रस्तावित माप विनिर्देश बनाते हैं, न कि यह कथन कि वर्तमान ऐप या प्रोटोकॉल प्रत्येक आवश्यक घटना को रिकॉर्ड करता है।

प्रत्येक प्रकाशित KPI में निम्नलिखित शामिल होंगेः

\begin{itemize}[leftmargin=*]
\item जिम्मेदार प्रकाशक और डेटा स्रोत;
\item घटना या राज्य की परिभाषा;
\item कोहर्ट या माप अवधि;
\item इकाई और मूल्य निर्धारण विधि;
\item मूल्यांकन समयशीर्षक;
\item समावेशन और बहिष्करण के नियम;
\item आंशिक, विवादास्पद, समाप्त हो चुके, अनुपलब्ध और सुधारित रिकॉर्डों का उपचार;
\item आवश्यक बाहर श्रृंखला सबूत या प्रमाणीकरण;
\item समीक्षा इतिहास और डेटा गुणवत्ता की सीमाएं; तथा
\item क्या परिणाम श्रृंखला पर है, रिपोर्ट किया गया, प्रमाणित, स्वतंत्र रूप से सत्यापित या अनुमानित।
\end{itemize}
\begin{longtable}{P{0.31\linewidth} P{0.31\linewidth} P{0.31\linewidth}}
\toprule
KPI & प्रस्तावित परिभाषा & आवश्यक साक्ष्य और बहिष्करण \\
\midrule
\endhead
\textbf{वैध प्रस्तुतियाँ} & एक अवधि के दौरान जारीकर्ता की प्रतिपूर्ति प्रक्रिया में स्वीकार किए गए अंक या इकाइयां & प्रस्तुति पहचान पत्र, जारीकर्ता, धारक प्राधिकरण, राशि, समय, स्थिति; डुप्लिकेट और अमान्य अनुरोधों को छोड़कर \\
\textbf{वादे की पूर्ति दर} & एक ही परिपक्व समूह के लिए वैध प्रस्तुतियों द्वारा विभाजित पूर्ण मान्य प्रस्तुतियाँ & अलग-अलग जारीकर्ता प्रदर्शन प्रमाण; खुले, अस्वीकृत, विवादास्पद, आंशिक और सुधारित मामलों की रिपोर्ट \\
\textbf{डिस्चार्ज पूर्णता} & पूर्ण प्रस्तुतिओं द्वारा विभाजित डिस्चार्ज रिकॉर्ड के साथ पूर्ण प्रस्तुतीकरण & जलाना, रद्द करना, निष्क्रिय करने वाला रिकॉर्ड या वादे की पूर्ति से जुड़ा अन्य गैर-पुनः उपयोग प्रमाण \\
\textbf{निष्पादन विलंबता} & प्रस्तुति-वादे की पूर्ति समय का मध्यवर्ती और 90वां प्रतिशत & जारी करने से प्रस्तुत करने या अधिग्रहण से प्रस्तुत होने तक की प्रतिधारण अवधि को प्रतिस्थापित न करें \\
\textbf{रखरखाव की अवधि} & अधिग्रहण से प्रस्तुत समय का औसत और वितरण & अधिग्रहण घटना की पहचान करें और अज्ञात अधिग्रहण समय को बाहर रखें \\
\textbf{शेष पात्र प्रतिबद्धताएं} & निर्दिष्ट बहिष्करणों के बाद शेष पात्र तृतीय-पक्ष प्रतिबद्धताएं & जारीकर्ता की पहचान और शर्तें; लागू जारीकर्ता इन्वेंट्री, समाप्ति, बर्न, डिस्चार्ज, परीक्षण और गैर-प्रतिबंध टोकन को छोड़कर \\
\textbf{पूल स्वैप की मात्रा} & एक खुलासा मूल्य निर्धारण पद्धति के तहत पूर्ण प्रत्यक्ष पूल स्वैप का मूल्य & ऑन-चेन सेटलमेंट इवेंट्स, टोकन इकाइयां, दर स्रोत, समय; जारीकर्ता वादे की पूर्ति के रूप में वर्गीकृत न करें \\
\textbf{पूल इन्वेंट्री} & निर्दिष्ट समय पर एक पूल द्वारा रखी गई मापी गई समर्थित परिसंपत्तियां & अनुबंध शेष राशि, शुल्क आरक्षण, inaccessible assets, valuation method और owner-withdrawal powers \\
\textbf{रिजर्व पर्याप्तता} & स्पष्ट रूप से कवर किए गए जोखिम द्वारा विभाजित पात्र, उपलब्ध आरक्षित परिसंपत्तियां & कवरेज नीति, परिसंपत्तियों की पात्रता, रखरखाव/नियंत्रण, दायित्व, बहिष्करण और मूल्यांकन; डिफ़ॉल्ट रूप से कुल टोकन आपूर्ति नहीं \\
\textbf{सीमा उपयोग} & वर्तमान कॉन्फ़िगर की गई कैप द्वारा विभाजित मापा गया पूल टोकन बैलेंस & Limiter पता, टोकन, पूल, टाइमस्टैम्प, परिवर्तन और अवधि बिना सीमांकित \\
\textbf{उद्धरण पास दर} & वैध उद्धरण प्रयासों द्वारा विभाजित सफल उद्धरण प्रतिक्रियाएं & केवल उद्धरण परिणाम; मार्ग निष्पादन नहीं \\
\textbf{मार्ग निष्पादन दर} & वैध निष्पादन प्रयासों द्वारा विभाजित पूर्ण मल्टी-हॉप निष्पादन & केवल एक कार्यान्वित निष्पादन प्रणाली पर लागू; प्रति-हॉप रिपोर्ट और परमाणुता नियम \\
\textbf{गारंटर की वसूली} & भुगतान किए गए कवर किए गए दावों में विभाजित प्राप्त पात्र वसूली & गारंटर की पहचान, दावा नीति, समय सीमा, लागत, विवाद और कटौती \\
\textbf{प्रस्तावित नेटवर्क राजस्व} & प्रस्तावित नेटवर्क रैक प्राप्त प्लस अलग-अलग रूटिंग/सेवा शुल्क प्राप्त & पूल द्वारा प्रतिधारण किए गए सकल पूल शुल्क को बाहर रखें और रैक को दो बार गिनने से बचें \\
\textbf{शासन की समयसीमा} & घटना का पता लगाने, निर्णय लेने, रुकने, मरम्मत करने और बंद करने का समय & निर्धारित घड़ियों, जिम्मेदार निकायों, आपातकालीन शक्तियों, अपीलों और गायब होने वाली घटनाओं \\
\bottomrule
\end{longtable}

सामाजिक प्रभाव के दावे को एक अलग पद्धति की आवश्यकता होती है। ब्लॉकचेन गतिविधि अकेले पहचान, जारीकर्ता प्रदर्शन, संतुष्टि, सामुदायिक स्वास्थ्य, कारणता या प्रभाव स्थापित नहीं करती है।


\section*{D. प्रस्तावित प्रक्षेपण मापदंड}
\addcontentsline{toc}{section}{D. प्रस्तावित प्रक्षेपण मापदंड}

नीचे दिए गए प्रत्येक मान चित्रणात्मक हैं. यह एक वर्तमान ऐप डिफ़ॉल्ट, Protocol v1.1.0 सेटिंग, गारंटी, पेशकश या वितरण प्रतिबद्धता नहीं है। भविष्य की तैनाती को अपने वास्तविक मापदंडों को अपनाने, लागू करने, निगरानी करने और प्रकाशित करने की आवश्यकता होगी।

\begin{itemize}[leftmargin=*]
\item \textbf{प्रस्तावित CLC शासन टोकन के लिए क्वोरम स्तरः}
\begin{itemize}[leftmargin=*]
\item Q1 रूटीनः कम से कम 4\% क्वोरम और 50\% से अधिक अनुमोदन।
\item Q2 संवेदनशीलः कम से कम 10\% क्वेरो और कम से अधिक 60\% अनुमोदन।
\item Q3महत्वपूर्ण: कम से कम 20\% क्विरोम और कम से66.7\% अनुमोदन।
\end{itemize}
\item \textbf{प्रस्तावित समय सीमाः}
\begin{itemize}[leftmargin=*]
\item T1: 48 घंटेQ1कार्य।
\item T2: 7 दिन के लिएQ2कार्य।
\item T3: 30 दिनQ3कार्य।
\end{itemize}
\item \textbf{युग का क्रमशः} 7 मतदान, बजट प्रकाशन और प्रस्तावित sCLC विंडो सहित दिन।
\item \textbf{आपातकालीन विराम:} तत्काल पहचान की गई आपातकालीन प्राधिकरण के माध्यम से, 72 घंटे के बाद समाप्त होता है जब तक कि इसे अपनाई गई प्रक्रिया के तहत अनुमोदित नहीं किया जाता।
\item \textbf{प्रस्तावित नेटवर्क रैकः} पॉलिसी द्वारा सीमित भाग लेने वाले पूल शुल्क का 20 प्रतिशत।
\item \textbf{इलस्ट्रेटिव पूल शुल्क सीमाः} 0\%--20\%, प्रत्येक भाग लेने वाले समूह द्वारा प्रकाशित।
\item \textbf{रूटिंग या सेवा शुल्क की प्रस्तावित सीमाः} 20 प्रति मार्ग बीपीएस।
\item \textbf{प्रस्तावित नेटवर्क रैक दरः} \texttt{rake\_rate\_p = pool\_fee\_p $\times$ rake\_share\_p}.
\item \textbf{राजस्व पात्रताः} प्राप्त परिसंपत्तियों को नकदी पात्र के रूप में वर्गीकृत करें\texttt{E\_cash} या प्रकृति में\texttt{E\_kind} प्रकाशित विधि के अनुसार।
\item \textbf{रूपांतरण नियंत्रणः} परिसंपत्तियों और स्थान अनुमतियों, जिम्मेदार अधिकारियों, मूल्य स्रोतों, समय खिड़कियों, स्लाइडिंग सीमाओं, कैप और रिपोर्टिंग।
\item \textbf{बजट की अनुमतिः} एक प्रस्तावित शुल्क-उपयोग बजट केवल तब प्रकाशित करें जब स्वीकृत कवर रिजर्व और परिचालन लक्ष्यों को पूरा किया गया हो; बजट शून्य हो सकता है।
\item \textbf{जोखिम लेन:} लागू कानून के तहत स्पष्ट, सूचित ऑप्ट-इन किए बिना एक परिसंपत्ति वर्ग को किसी अन्य वर्ग के जोखिम से नहीं निकालना।
\item \textbf{प्रस्तावित रोलिंग और खाता सीमाएंः} अनुमोदित क्रॉस-क्लास गतिविधि से इनकार करना और अपनाई गई सीमाओं को लागू करना।
\item \textbf{वैकल्पिक कवरेज में कमीः} केवल जब पहले से मौजूद कवरेज शर्तें, आवश्यक सहमति और लागू कानून इसे अधिकृत करते हैं; यह स्वयं किसी जारीकर्ता के वाउचर प्रतिबद्धता या श्रृंखला पर शेष राशि को कम नहीं करता है।
\item \textbf{बाहरी तरलता नियंत्रण, यदि अलग से सक्षम किया गया होः} स्थान के दायरे, माप विधि, ट्रिगर, व्यय और मात्रा की सीमाएं, निष्पादन नियंत्रण, आपातकालीन स्टॉप और रसीदें प्रकाशित करें। कार्यक्रम को टोकन-मूल्य समर्थन के रूप में प्रस्तुत न करें।
\end{itemize}
वर्तमान Protocol v1.1.0 पूल शुल्क, अतिरिक्त प्रोटोकॉल शुल्क और पूल टोकन बैलेंस कैप तैनात अनुबंधों और कॉन्फ़िगरेशन द्वारा शासित रहते हैं, न कि इन प्रस्तावित मूल्यों से।


\section*{E. कामकाजी उदाहरण  प्रस्तावित नेटवर्क रैक}
\addcontentsline{toc}{section}{E. कामकाजी उदाहरण  प्रस्तावित नेटवर्क रैक}

यह उदाहरण प्रस्तावित रैक मॉडल को दर्शाता है। यह Protocol v1.1.0 द्वारा उपयोग की जाने वाली अतिरिक्त प्रोटोकॉल-शुल्क गणना नहीं है।

मान लीजिएः

\begin{itemize}[leftmargin=*]
\item एक भाग लेने वाला पूल 2.00\% पूल शुल्क लेता है;
\item प्रस्तावित नेटवर्क रैक को उस पूल शुल्क का 20\% प्राप्त होता है; तथा
\item प्राप्त रैक का 25\% लागत के बाद घोषित नकदी में उपयोग के लिए पात्र और परिवर्तनीय है।
\end{itemize}
रूट किए गए मूल्य पर प्रभावी रूप से प्रस्तावित नेटवर्क रैक दर हैः

\texttt{rake\_rate = 2.00\% $\times$ 20\% = 0.40\% = 40 bps}

अनुमानित 25\% पात्रता कारक के तहत नकदी-उपयोग योग्य भाग हैः

\texttt{cash\_usable\_rate = 40 bps $\times$ 25\% = 10 bps}

यदि एक प्रस्तावित नेटवर्क बजट में \$ 150,000 और कोई अन्य राजस्व की मासिक नकदी विनिमय आवश्यकता है, तो 10 बीपीएस नकदी-उपयोग योग्य दर पर आवश्यक स्पष्ट रूप से रूटेड मूल्य हैः

\texttt{required\_routed\_value = \$150,000 / 0.001 = \$150,000,000 per month}

इस गणना में पूल द्वारा रखे गए सकल पूल शुल्क शामिल नहीं हैं। नेटवर्क रैक में इन शुल्कों को जोड़ने से रैक स्रोत की संख्या दोगुनी होगी।

एक वास्तविक बजट विश्लेषण में यह भी प्रकाशित करना आवश्यक होगाः

\begin{itemize}[leftmargin=*]
\item वास्तविक रैक और सेवा शुल्क प्राप्तियां;
\item परिसंपत्ति पात्रता और रूपांतरण लागत;
\item पूल भागीदारी और मार्ग की सीमा;
\item आरक्षित और परिचालन लक्ष्यों को अपनाया;
\item नुकसान, विवाद, सुधार और अनुपलब्ध संपत्ति; तथा
\item संवेदनाशीलता के परिदृश्यों की बजाय वादा किए गए विकास या रिटर्न।
\end{itemize}
शुल्क या तरलता कार्यक्रम के लिए कोई भी प्रस्तावित बजट अपनी स्वीकृत प्राथमिकताओं के निचले भाग में रहेगा और शून्य हो सकता है।


\section*{F. डाटा रूम चेकलिस्ट}
\addcontentsline{toc}{section}{F. डाटा रूम चेकलिस्ट}

\begin{enumerate}[leftmargin=*]
\item \textbf{ऐतिहासिक प्रमाण Sarafu Network}
\begin{itemize}[leftmargin=*]
\item अभिलेखागार दिनांकित स्वैप वॉल्यूम, सक्रिय उपयोगकर्ता और परिसंपत्ति परिभाषाएं, घटनाओं, चुकौती प्रस्तुतियाँ, उम्र बढ़ने वाले वाउचर, पुष्टि पूर्णता, डिस्चार्ज, रखरखाव अवधि उपायों, अवधि, बहिष्करण और पद्धति वर्तमान Cosmo-Local Credit मापदंडों से अलग।
\item संरक्षित करें \href{https://dune.com/grassrootseconomics/sarafu-network}{ऐतिहासिक ड्यून एनालिटिक्स डैशबोर्ड}.
\end{itemize}
\item \textbf{प्रस्तावित शुल्क निष्पादन प्रमाण, यदि लागू किया गया हो}
\begin{itemize}[leftmargin=*]
\item यह प्रदर्शित करें कि कोई भी शुल्क हुक और रजिस्टर गेटिंग केवल इंटरफ़ेस द्वारा लागू किए जाने के बजाय निष्पादित लेनदेन पथ में एकीकृत है।
\item ऐसे परीक्षण प्रदान करें जो यह दर्शाएं कि लागू निष्पादन सेवा एक उछाल को अस्वीकार करती है जिसकी प्राप्ति स्वीकृत नीति के तहत शुल्क संग्रह का प्रमाण नहीं देती है।
\item परीक्षण वेक्टर, विफलता के मामले, पात्र परिसंपत्तियां और एक पहचाने गए रजिस्ट्री संस्करण से जुड़े अनुरूप पूल प्रकाशित करें।
\item लागू लिंक \href{https://software.grassecon.org}{सॉफ्टवेयर विवरण}.
\end{itemize}
\item \textbf{गारंटी और कवरेज पैक}
\begin{itemize}[leftmargin=*]
\item पात्रता, उत्तरदायी पक्षों, धनराशि, सीमाएं, जहां लागू हो प्रीमियम, ट्रिगर, सबूत, बहिष्करण, निर्णय प्राधिकरण और विवाद और अपील के उदाहरण प्रकाशित करें।
\item नमूना वाउचर और पूल प्रकटीकरण शामिल करें जो किसी भी स्पष्ट रूप से प्रस्तावित पूल सुरक्षा या तृतीय-पक्ष गारंटी से जारीकर्ता की जिम्मेदारी को अलग करते हैं।
\item संरक्षित करें \href{https://sarafu.network/reports}{ऐतिहासिक Sarafu Network रिपोर्ट संग्रह} और \href{https://youtu.be/5yGaP31bvs0?si=fZ-AMTEXsmVR8Ulv}{समीक्षा के दौरान ऐतिहासिक वर्ष की प्रस्तुति} वंशावली के प्रमाण के रूप में, वर्तमान CLC कवरेज या गोद लेने का प्रमाण नहीं।
\end{itemize}
\item \textbf{कानूनी और अनुपालन पैक}
\begin{itemize}[leftmargin=*]
\item दस्तावेज क्षेत्राधिकार रणनीति, वाउचर वर्ग, नकद समकक्ष नीति, KYC या प्रमाणन विकल्प, भू-सीमांकन, उपभोक्ता प्रकटीकरण, गलत बिक्री नियंत्रण और सामान्य पूल स्वैप और एक्सप्रेस क्रेडिट उत्पाद के बीच अंतर।
\item प्रभावी को जोड़ें \href{https://docs.cosmolocal.credit/hi/governance/terms}{Cosmo-Local Credit सेवा की शर्तें} और नियंत्रित रिकॉर्ड के माध्यम से प्रतिस्थापित संस्करणों को संरक्षित करें।
\end{itemize}
\item \textbf{शासन सुदृढीकरण}
\begin{itemize}[leftmargin=*]
\item हितों के संघर्ष और अस्वीकरण की प्रक्रियाओं, निर्णय और अपील रिकॉर्ड, प्रतिबंध, रजिस्टर हटाने की उचित प्रक्रिया, आपातकालीन शक्ति सीमाएं, और समय में लॉक विनिमय दर और सीमा परिवर्तन के उदाहरण प्रदान करें।
\end{itemize}
\item \textbf{स्रोत और आश्वासन पैक}
\begin{itemize}[leftmargin=*]
\item स्रोत और लाइसेंस सूची, तैनात पते और संस्करणों, उपलब्ध ऑडिट रिपोर्ट, अपरिवर्तनीयताएं, निर्माण उत्पत्ति, व्यवस्थापक शक्तियां, घटना संपर्क और निगरानी डैशबोर्ड प्रकाशित करें।
\item लागू लिंक \href{https://github.com/grassrootseconomics}{जमीनी अर्थव्यवस्था भंडार}.
\end{itemize}
\end{enumerate}

\end{document}
